%!TEX program = xelatex
\documentclass[11pt,a4paper]{article}
\usepackage[margin=2.2cm]{geometry}
\usepackage{amsmath,amssymb,amsthm,mathtools}
\usepackage{bm}
\usepackage{enumitem}
\usepackage{booktabs}
\usepackage{array}
\usepackage{xcolor}
\usepackage{tikz}
\usetikzlibrary{arrows.meta,positioning}
\usepackage{hyperref}
\usepackage{fancyhdr}
\usepackage{xeCJK}

\hypersetup{
    colorlinks=true,
    linkcolor=blue!60!black,
    citecolor=blue!60!black,
    urlcolor=blue!60!black
}

\pagestyle{fancy}
\fancyhf{}
\rhead{\small Qiuzhen QE Computational \& Applied Math --- Spring 2025}
\lhead{\small Official Qualifying Exam}
\rfoot{\small Page \thepage}

\DeclareMathOperator*{\argmax}{arg\,max}
\DeclareMathOperator*{\argmin}{arg\,min}

\setlist[itemize]{topsep=3pt,itemsep=2pt,parsep=0pt}
\setlist[enumerate]{topsep=3pt,itemsep=2pt,parsep=0pt}

\newcommand{\examstatement}[1]{%
  \par\addvspace{0.85em}%
  \noindent{\bfseries #1}\par\nobreak\addvspace{0.28em}%
}

% Shared amsthm note/remark/theorem styles
% Shared math extras for qe-review.
% Requires already-loaded: amsmath, amssymb.
%
% Classic pitfall: AMS blackboard-bold fonts have no digit glyphs, so
% \mathbb{1} renders as overlapping garbage. Map the digit 1 to dsfont's
% \mathds{1}; leave \mathbb{R}, \mathbb{P}, etc. on AMS.
%
% Usage (path relative to the main .tex being compiled):
%   notes:  % Shared math extras for qe-review.
% Requires already-loaded: amsmath, amssymb.
%
% Classic pitfall: AMS blackboard-bold fonts have no digit glyphs, so
% \mathbb{1} renders as overlapping garbage. Map the digit 1 to dsfont's
% \mathds{1}; leave \mathbb{R}, \mathbb{P}, etc. on AMS.
%
% Usage (path relative to the main .tex being compiled):
%   notes:  % Shared math extras for qe-review.
% Requires already-loaded: amsmath, amssymb.
%
% Classic pitfall: AMS blackboard-bold fonts have no digit glyphs, so
% \mathbb{1} renders as overlapping garbage. Map the digit 1 to dsfont's
% \mathds{1}; leave \mathbb{R}, \mathbb{P}, etc. on AMS.
%
% Usage (path relative to the main .tex being compiled):
%   notes:  \input{../../common/qe-math-extras.tex}
%   exams:  \input{../../../common/qe-math-extras.tex}

\usepackage{dsfont}
\usepackage{pdftexcmds}

\makeatletter
\let\qe@amsmmathbb\mathbb
\renewcommand{\mathbb}[1]{%
  \ifnum\pdf@strcmp{\unexpanded{#1}}{1}=\z@
    \mathds{1}%
  \else
    \qe@amsmmathbb{#1}%
  \fi
}
\makeatother

% Preferred indicator macros (either style is safe after the remap above).
\newcommand{\bbone}{\mathds{1}}
\newcommand{\ind}{\mathbf{1}}

%   exams:  % Shared math extras for qe-review.
% Requires already-loaded: amsmath, amssymb.
%
% Classic pitfall: AMS blackboard-bold fonts have no digit glyphs, so
% \mathbb{1} renders as overlapping garbage. Map the digit 1 to dsfont's
% \mathds{1}; leave \mathbb{R}, \mathbb{P}, etc. on AMS.
%
% Usage (path relative to the main .tex being compiled):
%   notes:  \input{../../common/qe-math-extras.tex}
%   exams:  \input{../../../common/qe-math-extras.tex}

\usepackage{dsfont}
\usepackage{pdftexcmds}

\makeatletter
\let\qe@amsmmathbb\mathbb
\renewcommand{\mathbb}[1]{%
  \ifnum\pdf@strcmp{\unexpanded{#1}}{1}=\z@
    \mathds{1}%
  \else
    \qe@amsmmathbb{#1}%
  \fi
}
\makeatother

% Preferred indicator macros (either style is safe after the remap above).
\newcommand{\bbone}{\mathds{1}}
\newcommand{\ind}{\mathbf{1}}


\usepackage{dsfont}
\usepackage{pdftexcmds}

\makeatletter
\let\qe@amsmmathbb\mathbb
\renewcommand{\mathbb}[1]{%
  \ifnum\pdf@strcmp{\unexpanded{#1}}{1}=\z@
    \mathds{1}%
  \else
    \qe@amsmmathbb{#1}%
  \fi
}
\makeatother

% Preferred indicator macros (either style is safe after the remap above).
\newcommand{\bbone}{\mathds{1}}
\newcommand{\ind}{\mathbf{1}}

%   exams:  % Shared math extras for qe-review.
% Requires already-loaded: amsmath, amssymb.
%
% Classic pitfall: AMS blackboard-bold fonts have no digit glyphs, so
% \mathbb{1} renders as overlapping garbage. Map the digit 1 to dsfont's
% \mathds{1}; leave \mathbb{R}, \mathbb{P}, etc. on AMS.
%
% Usage (path relative to the main .tex being compiled):
%   notes:  % Shared math extras for qe-review.
% Requires already-loaded: amsmath, amssymb.
%
% Classic pitfall: AMS blackboard-bold fonts have no digit glyphs, so
% \mathbb{1} renders as overlapping garbage. Map the digit 1 to dsfont's
% \mathds{1}; leave \mathbb{R}, \mathbb{P}, etc. on AMS.
%
% Usage (path relative to the main .tex being compiled):
%   notes:  \input{../../common/qe-math-extras.tex}
%   exams:  \input{../../../common/qe-math-extras.tex}

\usepackage{dsfont}
\usepackage{pdftexcmds}

\makeatletter
\let\qe@amsmmathbb\mathbb
\renewcommand{\mathbb}[1]{%
  \ifnum\pdf@strcmp{\unexpanded{#1}}{1}=\z@
    \mathds{1}%
  \else
    \qe@amsmmathbb{#1}%
  \fi
}
\makeatother

% Preferred indicator macros (either style is safe after the remap above).
\newcommand{\bbone}{\mathds{1}}
\newcommand{\ind}{\mathbf{1}}

%   exams:  % Shared math extras for qe-review.
% Requires already-loaded: amsmath, amssymb.
%
% Classic pitfall: AMS blackboard-bold fonts have no digit glyphs, so
% \mathbb{1} renders as overlapping garbage. Map the digit 1 to dsfont's
% \mathds{1}; leave \mathbb{R}, \mathbb{P}, etc. on AMS.
%
% Usage (path relative to the main .tex being compiled):
%   notes:  \input{../../common/qe-math-extras.tex}
%   exams:  \input{../../../common/qe-math-extras.tex}

\usepackage{dsfont}
\usepackage{pdftexcmds}

\makeatletter
\let\qe@amsmmathbb\mathbb
\renewcommand{\mathbb}[1]{%
  \ifnum\pdf@strcmp{\unexpanded{#1}}{1}=\z@
    \mathds{1}%
  \else
    \qe@amsmmathbb{#1}%
  \fi
}
\makeatother

% Preferred indicator macros (either style is safe after the remap above).
\newcommand{\bbone}{\mathds{1}}
\newcommand{\ind}{\mathbf{1}}


\usepackage{dsfont}
\usepackage{pdftexcmds}

\makeatletter
\let\qe@amsmmathbb\mathbb
\renewcommand{\mathbb}[1]{%
  \ifnum\pdf@strcmp{\unexpanded{#1}}{1}=\z@
    \mathds{1}%
  \else
    \qe@amsmmathbb{#1}%
  \fi
}
\makeatother

% Preferred indicator macros (either style is safe after the remap above).
\newcommand{\bbone}{\mathds{1}}
\newcommand{\ind}{\mathbf{1}}


\usepackage{dsfont}
\usepackage{pdftexcmds}

\makeatletter
\let\qe@amsmmathbb\mathbb
\renewcommand{\mathbb}[1]{%
  \ifnum\pdf@strcmp{\unexpanded{#1}}{1}=\z@
    \mathds{1}%
  \else
    \qe@amsmmathbb{#1}%
  \fi
}
\makeatother

% Preferred indicator macros (either style is safe after the remap above).
\newcommand{\bbone}{\mathds{1}}
\newcommand{\ind}{\mathbf{1}}

% Shared amsthm environments for qe-review (minimal).
% Requires already-loaded: amsthm.
% Safe to \input after \usepackage{amsmath,amssymb,amsthm,...}.
%
% Shared numbering uses aliascnt so cleveref keys off the environment
% name (Lemma, Proposition, Exercise, ...) rather than the parent
% counter (Theorem, Definition, Remark). Without aliascnt, \cref{lem:...}
% prints ``Theorem 9.13'' instead of ``Lemma 9.13''.

\usepackage{aliascnt}

\theoremstyle{plain}
\newtheorem{theorem}{Theorem}[section]

\newaliascnt{lemma}{theorem}
\newtheorem{lemma}[lemma]{Lemma}
\aliascntresetthe{lemma}

\newaliascnt{proposition}{theorem}
\newtheorem{proposition}[proposition]{Proposition}
\aliascntresetthe{proposition}

\newaliascnt{corollary}{theorem}
\newtheorem{corollary}[corollary]{Corollary}
\aliascntresetthe{corollary}

\newaliascnt{claim}{theorem}
\newtheorem{claim}[claim]{Claim}
\aliascntresetthe{claim}

\newaliascnt{fact}{theorem}
\newtheorem{fact}[fact]{Fact}
\aliascntresetthe{fact}

\theoremstyle{definition}
\newtheorem{definition}{Definition}[section]

\newaliascnt{exercise}{definition}
\newtheorem{exercise}[exercise]{Exercise}
\aliascntresetthe{exercise}

\newaliascnt{assumption}{definition}
\newtheorem{assumption}[assumption]{Assumption}
\aliascntresetthe{assumption}

\newaliascnt{notation}{definition}
\newtheorem{notation}[notation]{Notation}
\aliascntresetthe{notation}

\newaliascnt{construction}{definition}
\newtheorem{construction}[construction]{Construction}
\aliascntresetthe{construction}

% Example: document-wide numbering (not per section / chapter)
\newtheorem{example}{Example}

\theoremstyle{remark}
\newtheorem{remark}{Remark}[section]
\newtheorem*{remark*}{Remark}

\newaliascnt{heuristic}{remark}
\newtheorem{heuristic}[heuristic]{Heuristic}
\aliascntresetthe{heuristic}

\newaliascnt{convention}{remark}
\newtheorem{convention}[convention]{Convention}
\aliascntresetthe{convention}

\newaliascnt{warning}{remark}
\newtheorem{warning}[warning]{Warning}
\aliascntresetthe{warning}

\newaliascnt{proofsketch}{remark}
\newtheorem{proofsketch}[proofsketch]{Proof sketch}
\aliascntresetthe{proofsketch}

\newaliascnt{checkpoint}{remark}
\newtheorem{checkpoint}[checkpoint]{Checkpoint}
\aliascntresetthe{checkpoint}

% Note: document-wide numbering (not per section)
\newtheorem{note}{Note}
\newtheorem*{note*}{Note}

\newaliascnt{examnote}{note}
\newtheorem{examnote}[examnote]{Exam note}
\aliascntresetthe{examnote}

\newaliascnt{study}{note}
\newtheorem{study}[study]{Study note}
\aliascntresetthe{study}

\newaliascnt{summary}{note}
\newtheorem{summary}[summary]{Summary}
\aliascntresetthe{summary}

\newtheorem{examproblem}{Problem}[section]
\newtheorem*{examproblem*}{Problem}

% Bilingual aliases
\newenvironment{dingli}{\begin{theorem}}{\end{theorem}}
\newenvironment{yinli}{\begin{lemma}}{\end{lemma}}
\newenvironment{mingti}{\begin{proposition}}{\end{proposition}}
\newenvironment{tuilun}{\begin{corollary}}{\end{corollary}}
\newenvironment{dingyi}{\begin{definition}}{\end{definition}}
\newenvironment{li}{\begin{example}}{\end{example}}
\newenvironment{zhu}{\begin{remark}}{\end{remark}}
\newenvironment{biji}{\begin{note}}{\end{note}}
\newenvironment{kaoshi}{\begin{examnote}}{\end{examnote}}

\renewcommand{\qedsymbol}{\ensuremath{\square}}

% Shared cleveref setup for qe-review.
% Load AFTER: hyperref, and AFTER all \newtheorem / amsthm environments.
% Shared theorem counters are aliased in qe-amsthm-envs.tex (aliascnt),
% otherwise \cref on a lemma/proposition prints ``Theorem''.
%
% Usage (path relative to the main .tex being compiled):
%   notes:  % Shared cleveref setup for qe-review.
% Load AFTER: hyperref, and AFTER all \newtheorem / amsthm environments.
% Shared theorem counters are aliased in qe-amsthm-envs.tex (aliascnt),
% otherwise \cref on a lemma/proposition prints ``Theorem''.
%
% Usage (path relative to the main .tex being compiled):
%   notes:  % Shared cleveref setup for qe-review.
% Load AFTER: hyperref, and AFTER all \newtheorem / amsthm environments.
% Shared theorem counters are aliased in qe-amsthm-envs.tex (aliascnt),
% otherwise \cref on a lemma/proposition prints ``Theorem''.
%
% Usage (path relative to the main .tex being compiled):
%   notes:  \input{../../common/qe-cleveref.tex}
%   exams:  \input{../../../common/qe-cleveref.tex}
%
% Prefer \cref / \Cref over \ref. Use \Cref at sentence start.
% Unnumbered sectioning (e.g. \paragraph with default secnumdepth)
% produces an empty cleveref type; map that to "Paragraph".

\usepackage[nameinlink,noabbrev]{cleveref}

% Fallback for unnumbered \paragraph / \subsubsection labels
\crefname{}{Paragraph}{Paragraphs}
\Crefname{}{Paragraph}{Paragraphs}

% Numbered theorem-like (plain)
\crefname{theorem}{Theorem}{Theorems}
\Crefname{theorem}{Theorem}{Theorems}
\crefname{lemma}{Lemma}{Lemmas}
\Crefname{lemma}{Lemma}{Lemmas}
\crefname{proposition}{Proposition}{Propositions}
\Crefname{proposition}{Proposition}{Propositions}
\crefname{corollary}{Corollary}{Corollaries}
\Crefname{corollary}{Corollary}{Corollaries}
\crefname{claim}{Claim}{Claims}
\Crefname{claim}{Claim}{Claims}
\crefname{fact}{Fact}{Facts}
\Crefname{fact}{Fact}{Facts}

% Definition-like
\crefname{definition}{Definition}{Definitions}
\Crefname{definition}{Definition}{Definitions}
\crefname{example}{Example}{Examples}
\Crefname{example}{Example}{Examples}
\crefname{exercise}{Exercise}{Exercises}
\Crefname{exercise}{Exercise}{Exercises}
\crefname{assumption}{Assumption}{Assumptions}
\Crefname{assumption}{Assumption}{Assumptions}
\crefname{notation}{Notation}{Notations}
\Crefname{notation}{Notation}{Notations}
\crefname{construction}{Construction}{Constructions}
\Crefname{construction}{Construction}{Constructions}

% Remark-like
\crefname{remark}{Remark}{Remarks}
\Crefname{remark}{Remark}{Remarks}
\crefname{heuristic}{Heuristic}{Heuristics}
\Crefname{heuristic}{Heuristic}{Heuristics}
\crefname{convention}{Convention}{Conventions}
\Crefname{convention}{Convention}{Conventions}
\crefname{warning}{Warning}{Warnings}
\Crefname{warning}{Warning}{Warnings}
\crefname{proofsketch}{Proof sketch}{Proof sketches}
\Crefname{proofsketch}{Proof sketch}{Proof sketches}
\crefname{checkpoint}{Checkpoint}{Checkpoints}
\Crefname{checkpoint}{Checkpoint}{Checkpoints}

% Document-wide notes / exam problems
\crefname{note}{Note}{Notes}
\Crefname{note}{Note}{Notes}
\crefname{examnote}{Exam note}{Exam notes}
\Crefname{examnote}{Exam note}{Exam notes}
\crefname{study}{Study note}{Study notes}
\Crefname{study}{Study note}{Study notes}
\crefname{summary}{Summary}{Summaries}
\Crefname{summary}{Summary}{Summaries}
\crefname{examproblem}{Problem}{Problems}
\Crefname{examproblem}{Problem}{Problems}

% Sectioning / floats (explicit, noabbrev-friendly)
\crefname{section}{Section}{Sections}
\Crefname{section}{Section}{Sections}
\crefname{subsection}{Subsection}{Subsections}
\Crefname{subsection}{Subsection}{Subsections}
\crefname{subsubsection}{Subsubsection}{Subsubsections}
\Crefname{subsubsection}{Subsubsection}{Subsubsections}
\crefname{paragraph}{Paragraph}{Paragraphs}
\Crefname{paragraph}{Paragraph}{Paragraphs}
\crefname{equation}{Equation}{Equations}
\Crefname{equation}{Equation}{Equations}
\crefname{figure}{Figure}{Figures}
\Crefname{figure}{Figure}{Figures}
\crefname{table}{Table}{Tables}
\Crefname{table}{Table}{Tables}
\crefname{enumi}{Item}{Items}
\Crefname{enumi}{Item}{Items}

%   exams:  % Shared cleveref setup for qe-review.
% Load AFTER: hyperref, and AFTER all \newtheorem / amsthm environments.
% Shared theorem counters are aliased in qe-amsthm-envs.tex (aliascnt),
% otherwise \cref on a lemma/proposition prints ``Theorem''.
%
% Usage (path relative to the main .tex being compiled):
%   notes:  \input{../../common/qe-cleveref.tex}
%   exams:  \input{../../../common/qe-cleveref.tex}
%
% Prefer \cref / \Cref over \ref. Use \Cref at sentence start.
% Unnumbered sectioning (e.g. \paragraph with default secnumdepth)
% produces an empty cleveref type; map that to "Paragraph".

\usepackage[nameinlink,noabbrev]{cleveref}

% Fallback for unnumbered \paragraph / \subsubsection labels
\crefname{}{Paragraph}{Paragraphs}
\Crefname{}{Paragraph}{Paragraphs}

% Numbered theorem-like (plain)
\crefname{theorem}{Theorem}{Theorems}
\Crefname{theorem}{Theorem}{Theorems}
\crefname{lemma}{Lemma}{Lemmas}
\Crefname{lemma}{Lemma}{Lemmas}
\crefname{proposition}{Proposition}{Propositions}
\Crefname{proposition}{Proposition}{Propositions}
\crefname{corollary}{Corollary}{Corollaries}
\Crefname{corollary}{Corollary}{Corollaries}
\crefname{claim}{Claim}{Claims}
\Crefname{claim}{Claim}{Claims}
\crefname{fact}{Fact}{Facts}
\Crefname{fact}{Fact}{Facts}

% Definition-like
\crefname{definition}{Definition}{Definitions}
\Crefname{definition}{Definition}{Definitions}
\crefname{example}{Example}{Examples}
\Crefname{example}{Example}{Examples}
\crefname{exercise}{Exercise}{Exercises}
\Crefname{exercise}{Exercise}{Exercises}
\crefname{assumption}{Assumption}{Assumptions}
\Crefname{assumption}{Assumption}{Assumptions}
\crefname{notation}{Notation}{Notations}
\Crefname{notation}{Notation}{Notations}
\crefname{construction}{Construction}{Constructions}
\Crefname{construction}{Construction}{Constructions}

% Remark-like
\crefname{remark}{Remark}{Remarks}
\Crefname{remark}{Remark}{Remarks}
\crefname{heuristic}{Heuristic}{Heuristics}
\Crefname{heuristic}{Heuristic}{Heuristics}
\crefname{convention}{Convention}{Conventions}
\Crefname{convention}{Convention}{Conventions}
\crefname{warning}{Warning}{Warnings}
\Crefname{warning}{Warning}{Warnings}
\crefname{proofsketch}{Proof sketch}{Proof sketches}
\Crefname{proofsketch}{Proof sketch}{Proof sketches}
\crefname{checkpoint}{Checkpoint}{Checkpoints}
\Crefname{checkpoint}{Checkpoint}{Checkpoints}

% Document-wide notes / exam problems
\crefname{note}{Note}{Notes}
\Crefname{note}{Note}{Notes}
\crefname{examnote}{Exam note}{Exam notes}
\Crefname{examnote}{Exam note}{Exam notes}
\crefname{study}{Study note}{Study notes}
\Crefname{study}{Study note}{Study notes}
\crefname{summary}{Summary}{Summaries}
\Crefname{summary}{Summary}{Summaries}
\crefname{examproblem}{Problem}{Problems}
\Crefname{examproblem}{Problem}{Problems}

% Sectioning / floats (explicit, noabbrev-friendly)
\crefname{section}{Section}{Sections}
\Crefname{section}{Section}{Sections}
\crefname{subsection}{Subsection}{Subsections}
\Crefname{subsection}{Subsection}{Subsections}
\crefname{subsubsection}{Subsubsection}{Subsubsections}
\Crefname{subsubsection}{Subsubsection}{Subsubsections}
\crefname{paragraph}{Paragraph}{Paragraphs}
\Crefname{paragraph}{Paragraph}{Paragraphs}
\crefname{equation}{Equation}{Equations}
\Crefname{equation}{Equation}{Equations}
\crefname{figure}{Figure}{Figures}
\Crefname{figure}{Figure}{Figures}
\crefname{table}{Table}{Tables}
\Crefname{table}{Table}{Tables}
\crefname{enumi}{Item}{Items}
\Crefname{enumi}{Item}{Items}

%
% Prefer \cref / \Cref over \ref. Use \Cref at sentence start.
% Unnumbered sectioning (e.g. \paragraph with default secnumdepth)
% produces an empty cleveref type; map that to "Paragraph".

\usepackage[nameinlink,noabbrev]{cleveref}

% Fallback for unnumbered \paragraph / \subsubsection labels
\crefname{}{Paragraph}{Paragraphs}
\Crefname{}{Paragraph}{Paragraphs}

% Numbered theorem-like (plain)
\crefname{theorem}{Theorem}{Theorems}
\Crefname{theorem}{Theorem}{Theorems}
\crefname{lemma}{Lemma}{Lemmas}
\Crefname{lemma}{Lemma}{Lemmas}
\crefname{proposition}{Proposition}{Propositions}
\Crefname{proposition}{Proposition}{Propositions}
\crefname{corollary}{Corollary}{Corollaries}
\Crefname{corollary}{Corollary}{Corollaries}
\crefname{claim}{Claim}{Claims}
\Crefname{claim}{Claim}{Claims}
\crefname{fact}{Fact}{Facts}
\Crefname{fact}{Fact}{Facts}

% Definition-like
\crefname{definition}{Definition}{Definitions}
\Crefname{definition}{Definition}{Definitions}
\crefname{example}{Example}{Examples}
\Crefname{example}{Example}{Examples}
\crefname{exercise}{Exercise}{Exercises}
\Crefname{exercise}{Exercise}{Exercises}
\crefname{assumption}{Assumption}{Assumptions}
\Crefname{assumption}{Assumption}{Assumptions}
\crefname{notation}{Notation}{Notations}
\Crefname{notation}{Notation}{Notations}
\crefname{construction}{Construction}{Constructions}
\Crefname{construction}{Construction}{Constructions}

% Remark-like
\crefname{remark}{Remark}{Remarks}
\Crefname{remark}{Remark}{Remarks}
\crefname{heuristic}{Heuristic}{Heuristics}
\Crefname{heuristic}{Heuristic}{Heuristics}
\crefname{convention}{Convention}{Conventions}
\Crefname{convention}{Convention}{Conventions}
\crefname{warning}{Warning}{Warnings}
\Crefname{warning}{Warning}{Warnings}
\crefname{proofsketch}{Proof sketch}{Proof sketches}
\Crefname{proofsketch}{Proof sketch}{Proof sketches}
\crefname{checkpoint}{Checkpoint}{Checkpoints}
\Crefname{checkpoint}{Checkpoint}{Checkpoints}

% Document-wide notes / exam problems
\crefname{note}{Note}{Notes}
\Crefname{note}{Note}{Notes}
\crefname{examnote}{Exam note}{Exam notes}
\Crefname{examnote}{Exam note}{Exam notes}
\crefname{study}{Study note}{Study notes}
\Crefname{study}{Study note}{Study notes}
\crefname{summary}{Summary}{Summaries}
\Crefname{summary}{Summary}{Summaries}
\crefname{examproblem}{Problem}{Problems}
\Crefname{examproblem}{Problem}{Problems}

% Sectioning / floats (explicit, noabbrev-friendly)
\crefname{section}{Section}{Sections}
\Crefname{section}{Section}{Sections}
\crefname{subsection}{Subsection}{Subsections}
\Crefname{subsection}{Subsection}{Subsections}
\crefname{subsubsection}{Subsubsection}{Subsubsections}
\Crefname{subsubsection}{Subsubsection}{Subsubsections}
\crefname{paragraph}{Paragraph}{Paragraphs}
\Crefname{paragraph}{Paragraph}{Paragraphs}
\crefname{equation}{Equation}{Equations}
\Crefname{equation}{Equation}{Equations}
\crefname{figure}{Figure}{Figures}
\Crefname{figure}{Figure}{Figures}
\crefname{table}{Table}{Tables}
\Crefname{table}{Table}{Tables}
\crefname{enumi}{Item}{Items}
\Crefname{enumi}{Item}{Items}

%   exams:  % Shared cleveref setup for qe-review.
% Load AFTER: hyperref, and AFTER all \newtheorem / amsthm environments.
% Shared theorem counters are aliased in qe-amsthm-envs.tex (aliascnt),
% otherwise \cref on a lemma/proposition prints ``Theorem''.
%
% Usage (path relative to the main .tex being compiled):
%   notes:  % Shared cleveref setup for qe-review.
% Load AFTER: hyperref, and AFTER all \newtheorem / amsthm environments.
% Shared theorem counters are aliased in qe-amsthm-envs.tex (aliascnt),
% otherwise \cref on a lemma/proposition prints ``Theorem''.
%
% Usage (path relative to the main .tex being compiled):
%   notes:  \input{../../common/qe-cleveref.tex}
%   exams:  \input{../../../common/qe-cleveref.tex}
%
% Prefer \cref / \Cref over \ref. Use \Cref at sentence start.
% Unnumbered sectioning (e.g. \paragraph with default secnumdepth)
% produces an empty cleveref type; map that to "Paragraph".

\usepackage[nameinlink,noabbrev]{cleveref}

% Fallback for unnumbered \paragraph / \subsubsection labels
\crefname{}{Paragraph}{Paragraphs}
\Crefname{}{Paragraph}{Paragraphs}

% Numbered theorem-like (plain)
\crefname{theorem}{Theorem}{Theorems}
\Crefname{theorem}{Theorem}{Theorems}
\crefname{lemma}{Lemma}{Lemmas}
\Crefname{lemma}{Lemma}{Lemmas}
\crefname{proposition}{Proposition}{Propositions}
\Crefname{proposition}{Proposition}{Propositions}
\crefname{corollary}{Corollary}{Corollaries}
\Crefname{corollary}{Corollary}{Corollaries}
\crefname{claim}{Claim}{Claims}
\Crefname{claim}{Claim}{Claims}
\crefname{fact}{Fact}{Facts}
\Crefname{fact}{Fact}{Facts}

% Definition-like
\crefname{definition}{Definition}{Definitions}
\Crefname{definition}{Definition}{Definitions}
\crefname{example}{Example}{Examples}
\Crefname{example}{Example}{Examples}
\crefname{exercise}{Exercise}{Exercises}
\Crefname{exercise}{Exercise}{Exercises}
\crefname{assumption}{Assumption}{Assumptions}
\Crefname{assumption}{Assumption}{Assumptions}
\crefname{notation}{Notation}{Notations}
\Crefname{notation}{Notation}{Notations}
\crefname{construction}{Construction}{Constructions}
\Crefname{construction}{Construction}{Constructions}

% Remark-like
\crefname{remark}{Remark}{Remarks}
\Crefname{remark}{Remark}{Remarks}
\crefname{heuristic}{Heuristic}{Heuristics}
\Crefname{heuristic}{Heuristic}{Heuristics}
\crefname{convention}{Convention}{Conventions}
\Crefname{convention}{Convention}{Conventions}
\crefname{warning}{Warning}{Warnings}
\Crefname{warning}{Warning}{Warnings}
\crefname{proofsketch}{Proof sketch}{Proof sketches}
\Crefname{proofsketch}{Proof sketch}{Proof sketches}
\crefname{checkpoint}{Checkpoint}{Checkpoints}
\Crefname{checkpoint}{Checkpoint}{Checkpoints}

% Document-wide notes / exam problems
\crefname{note}{Note}{Notes}
\Crefname{note}{Note}{Notes}
\crefname{examnote}{Exam note}{Exam notes}
\Crefname{examnote}{Exam note}{Exam notes}
\crefname{study}{Study note}{Study notes}
\Crefname{study}{Study note}{Study notes}
\crefname{summary}{Summary}{Summaries}
\Crefname{summary}{Summary}{Summaries}
\crefname{examproblem}{Problem}{Problems}
\Crefname{examproblem}{Problem}{Problems}

% Sectioning / floats (explicit, noabbrev-friendly)
\crefname{section}{Section}{Sections}
\Crefname{section}{Section}{Sections}
\crefname{subsection}{Subsection}{Subsections}
\Crefname{subsection}{Subsection}{Subsections}
\crefname{subsubsection}{Subsubsection}{Subsubsections}
\Crefname{subsubsection}{Subsubsection}{Subsubsections}
\crefname{paragraph}{Paragraph}{Paragraphs}
\Crefname{paragraph}{Paragraph}{Paragraphs}
\crefname{equation}{Equation}{Equations}
\Crefname{equation}{Equation}{Equations}
\crefname{figure}{Figure}{Figures}
\Crefname{figure}{Figure}{Figures}
\crefname{table}{Table}{Tables}
\Crefname{table}{Table}{Tables}
\crefname{enumi}{Item}{Items}
\Crefname{enumi}{Item}{Items}

%   exams:  % Shared cleveref setup for qe-review.
% Load AFTER: hyperref, and AFTER all \newtheorem / amsthm environments.
% Shared theorem counters are aliased in qe-amsthm-envs.tex (aliascnt),
% otherwise \cref on a lemma/proposition prints ``Theorem''.
%
% Usage (path relative to the main .tex being compiled):
%   notes:  \input{../../common/qe-cleveref.tex}
%   exams:  \input{../../../common/qe-cleveref.tex}
%
% Prefer \cref / \Cref over \ref. Use \Cref at sentence start.
% Unnumbered sectioning (e.g. \paragraph with default secnumdepth)
% produces an empty cleveref type; map that to "Paragraph".

\usepackage[nameinlink,noabbrev]{cleveref}

% Fallback for unnumbered \paragraph / \subsubsection labels
\crefname{}{Paragraph}{Paragraphs}
\Crefname{}{Paragraph}{Paragraphs}

% Numbered theorem-like (plain)
\crefname{theorem}{Theorem}{Theorems}
\Crefname{theorem}{Theorem}{Theorems}
\crefname{lemma}{Lemma}{Lemmas}
\Crefname{lemma}{Lemma}{Lemmas}
\crefname{proposition}{Proposition}{Propositions}
\Crefname{proposition}{Proposition}{Propositions}
\crefname{corollary}{Corollary}{Corollaries}
\Crefname{corollary}{Corollary}{Corollaries}
\crefname{claim}{Claim}{Claims}
\Crefname{claim}{Claim}{Claims}
\crefname{fact}{Fact}{Facts}
\Crefname{fact}{Fact}{Facts}

% Definition-like
\crefname{definition}{Definition}{Definitions}
\Crefname{definition}{Definition}{Definitions}
\crefname{example}{Example}{Examples}
\Crefname{example}{Example}{Examples}
\crefname{exercise}{Exercise}{Exercises}
\Crefname{exercise}{Exercise}{Exercises}
\crefname{assumption}{Assumption}{Assumptions}
\Crefname{assumption}{Assumption}{Assumptions}
\crefname{notation}{Notation}{Notations}
\Crefname{notation}{Notation}{Notations}
\crefname{construction}{Construction}{Constructions}
\Crefname{construction}{Construction}{Constructions}

% Remark-like
\crefname{remark}{Remark}{Remarks}
\Crefname{remark}{Remark}{Remarks}
\crefname{heuristic}{Heuristic}{Heuristics}
\Crefname{heuristic}{Heuristic}{Heuristics}
\crefname{convention}{Convention}{Conventions}
\Crefname{convention}{Convention}{Conventions}
\crefname{warning}{Warning}{Warnings}
\Crefname{warning}{Warning}{Warnings}
\crefname{proofsketch}{Proof sketch}{Proof sketches}
\Crefname{proofsketch}{Proof sketch}{Proof sketches}
\crefname{checkpoint}{Checkpoint}{Checkpoints}
\Crefname{checkpoint}{Checkpoint}{Checkpoints}

% Document-wide notes / exam problems
\crefname{note}{Note}{Notes}
\Crefname{note}{Note}{Notes}
\crefname{examnote}{Exam note}{Exam notes}
\Crefname{examnote}{Exam note}{Exam notes}
\crefname{study}{Study note}{Study notes}
\Crefname{study}{Study note}{Study notes}
\crefname{summary}{Summary}{Summaries}
\Crefname{summary}{Summary}{Summaries}
\crefname{examproblem}{Problem}{Problems}
\Crefname{examproblem}{Problem}{Problems}

% Sectioning / floats (explicit, noabbrev-friendly)
\crefname{section}{Section}{Sections}
\Crefname{section}{Section}{Sections}
\crefname{subsection}{Subsection}{Subsections}
\Crefname{subsection}{Subsection}{Subsections}
\crefname{subsubsection}{Subsubsection}{Subsubsections}
\Crefname{subsubsection}{Subsubsection}{Subsubsections}
\crefname{paragraph}{Paragraph}{Paragraphs}
\Crefname{paragraph}{Paragraph}{Paragraphs}
\crefname{equation}{Equation}{Equations}
\Crefname{equation}{Equation}{Equations}
\crefname{figure}{Figure}{Figures}
\Crefname{figure}{Figure}{Figures}
\crefname{table}{Table}{Tables}
\Crefname{table}{Table}{Tables}
\crefname{enumi}{Item}{Items}
\Crefname{enumi}{Item}{Items}

%
% Prefer \cref / \Cref over \ref. Use \Cref at sentence start.
% Unnumbered sectioning (e.g. \paragraph with default secnumdepth)
% produces an empty cleveref type; map that to "Paragraph".

\usepackage[nameinlink,noabbrev]{cleveref}

% Fallback for unnumbered \paragraph / \subsubsection labels
\crefname{}{Paragraph}{Paragraphs}
\Crefname{}{Paragraph}{Paragraphs}

% Numbered theorem-like (plain)
\crefname{theorem}{Theorem}{Theorems}
\Crefname{theorem}{Theorem}{Theorems}
\crefname{lemma}{Lemma}{Lemmas}
\Crefname{lemma}{Lemma}{Lemmas}
\crefname{proposition}{Proposition}{Propositions}
\Crefname{proposition}{Proposition}{Propositions}
\crefname{corollary}{Corollary}{Corollaries}
\Crefname{corollary}{Corollary}{Corollaries}
\crefname{claim}{Claim}{Claims}
\Crefname{claim}{Claim}{Claims}
\crefname{fact}{Fact}{Facts}
\Crefname{fact}{Fact}{Facts}

% Definition-like
\crefname{definition}{Definition}{Definitions}
\Crefname{definition}{Definition}{Definitions}
\crefname{example}{Example}{Examples}
\Crefname{example}{Example}{Examples}
\crefname{exercise}{Exercise}{Exercises}
\Crefname{exercise}{Exercise}{Exercises}
\crefname{assumption}{Assumption}{Assumptions}
\Crefname{assumption}{Assumption}{Assumptions}
\crefname{notation}{Notation}{Notations}
\Crefname{notation}{Notation}{Notations}
\crefname{construction}{Construction}{Constructions}
\Crefname{construction}{Construction}{Constructions}

% Remark-like
\crefname{remark}{Remark}{Remarks}
\Crefname{remark}{Remark}{Remarks}
\crefname{heuristic}{Heuristic}{Heuristics}
\Crefname{heuristic}{Heuristic}{Heuristics}
\crefname{convention}{Convention}{Conventions}
\Crefname{convention}{Convention}{Conventions}
\crefname{warning}{Warning}{Warnings}
\Crefname{warning}{Warning}{Warnings}
\crefname{proofsketch}{Proof sketch}{Proof sketches}
\Crefname{proofsketch}{Proof sketch}{Proof sketches}
\crefname{checkpoint}{Checkpoint}{Checkpoints}
\Crefname{checkpoint}{Checkpoint}{Checkpoints}

% Document-wide notes / exam problems
\crefname{note}{Note}{Notes}
\Crefname{note}{Note}{Notes}
\crefname{examnote}{Exam note}{Exam notes}
\Crefname{examnote}{Exam note}{Exam notes}
\crefname{study}{Study note}{Study notes}
\Crefname{study}{Study note}{Study notes}
\crefname{summary}{Summary}{Summaries}
\Crefname{summary}{Summary}{Summaries}
\crefname{examproblem}{Problem}{Problems}
\Crefname{examproblem}{Problem}{Problems}

% Sectioning / floats (explicit, noabbrev-friendly)
\crefname{section}{Section}{Sections}
\Crefname{section}{Section}{Sections}
\crefname{subsection}{Subsection}{Subsections}
\Crefname{subsection}{Subsection}{Subsections}
\crefname{subsubsection}{Subsubsection}{Subsubsections}
\Crefname{subsubsection}{Subsubsection}{Subsubsections}
\crefname{paragraph}{Paragraph}{Paragraphs}
\Crefname{paragraph}{Paragraph}{Paragraphs}
\crefname{equation}{Equation}{Equations}
\Crefname{equation}{Equation}{Equations}
\crefname{figure}{Figure}{Figures}
\Crefname{figure}{Figure}{Figures}
\crefname{table}{Table}{Tables}
\Crefname{table}{Table}{Tables}
\crefname{enumi}{Item}{Items}
\Crefname{enumi}{Item}{Items}

%
% Prefer \cref / \Cref over \ref. Use \Cref at sentence start.
% Unnumbered sectioning (e.g. \paragraph with default secnumdepth)
% produces an empty cleveref type; map that to "Paragraph".

\usepackage[nameinlink,noabbrev]{cleveref}

% Fallback for unnumbered \paragraph / \subsubsection labels
\crefname{}{Paragraph}{Paragraphs}
\Crefname{}{Paragraph}{Paragraphs}

% Numbered theorem-like (plain)
\crefname{theorem}{Theorem}{Theorems}
\Crefname{theorem}{Theorem}{Theorems}
\crefname{lemma}{Lemma}{Lemmas}
\Crefname{lemma}{Lemma}{Lemmas}
\crefname{proposition}{Proposition}{Propositions}
\Crefname{proposition}{Proposition}{Propositions}
\crefname{corollary}{Corollary}{Corollaries}
\Crefname{corollary}{Corollary}{Corollaries}
\crefname{claim}{Claim}{Claims}
\Crefname{claim}{Claim}{Claims}
\crefname{fact}{Fact}{Facts}
\Crefname{fact}{Fact}{Facts}

% Definition-like
\crefname{definition}{Definition}{Definitions}
\Crefname{definition}{Definition}{Definitions}
\crefname{example}{Example}{Examples}
\Crefname{example}{Example}{Examples}
\crefname{exercise}{Exercise}{Exercises}
\Crefname{exercise}{Exercise}{Exercises}
\crefname{assumption}{Assumption}{Assumptions}
\Crefname{assumption}{Assumption}{Assumptions}
\crefname{notation}{Notation}{Notations}
\Crefname{notation}{Notation}{Notations}
\crefname{construction}{Construction}{Constructions}
\Crefname{construction}{Construction}{Constructions}

% Remark-like
\crefname{remark}{Remark}{Remarks}
\Crefname{remark}{Remark}{Remarks}
\crefname{heuristic}{Heuristic}{Heuristics}
\Crefname{heuristic}{Heuristic}{Heuristics}
\crefname{convention}{Convention}{Conventions}
\Crefname{convention}{Convention}{Conventions}
\crefname{warning}{Warning}{Warnings}
\Crefname{warning}{Warning}{Warnings}
\crefname{proofsketch}{Proof sketch}{Proof sketches}
\Crefname{proofsketch}{Proof sketch}{Proof sketches}
\crefname{checkpoint}{Checkpoint}{Checkpoints}
\Crefname{checkpoint}{Checkpoint}{Checkpoints}

% Document-wide notes / exam problems
\crefname{note}{Note}{Notes}
\Crefname{note}{Note}{Notes}
\crefname{examnote}{Exam note}{Exam notes}
\Crefname{examnote}{Exam note}{Exam notes}
\crefname{study}{Study note}{Study notes}
\Crefname{study}{Study note}{Study notes}
\crefname{summary}{Summary}{Summaries}
\Crefname{summary}{Summary}{Summaries}
\crefname{examproblem}{Problem}{Problems}
\Crefname{examproblem}{Problem}{Problems}

% Sectioning / floats (explicit, noabbrev-friendly)
\crefname{section}{Section}{Sections}
\Crefname{section}{Section}{Sections}
\crefname{subsection}{Subsection}{Subsections}
\Crefname{subsection}{Subsection}{Subsections}
\crefname{subsubsection}{Subsubsection}{Subsubsections}
\Crefname{subsubsection}{Subsubsection}{Subsubsections}
\crefname{paragraph}{Paragraph}{Paragraphs}
\Crefname{paragraph}{Paragraph}{Paragraphs}
\crefname{equation}{Equation}{Equations}
\Crefname{equation}{Equation}{Equations}
\crefname{figure}{Figure}{Figures}
\Crefname{figure}{Figure}{Figures}
\crefname{table}{Table}{Tables}
\Crefname{table}{Table}{Tables}
\crefname{enumi}{Item}{Items}
\Crefname{enumi}{Item}{Items}

% PDF sidebar outline + printed TOC for QE exam transcripts.
% Load in the preamble AFTER hyperref and AFTER \newcommand{\examstatement}.
\hypersetup{
  unicode=true,
  bookmarks=true,
  bookmarksopen=true,
  bookmarksopenlevel=3,
  bookmarksnumbered=false,
  bookmarksdepth=3
}
\IfFileExists{bookmark.sty}{%
  \usepackage{bookmark}%
  \bookmarksetup{open,openlevel=3,numbered=false,depth=3}%
}{}
% Unnumbered in print, but still written to the TOC / PDF outline
% down to \subsubsection. Starred headings create neither.
\setcounter{secnumdepth}{-1}
\setcounter{tocdepth}{3}
\makeatletter
\@ifundefined{examstatement}{}{%
  \let\qe@origexamstatement\examstatement
  \renewcommand{\examstatement}[1]{%
    \par\phantomsection
    \addcontentsline{toc}{subsection}{#1}%
    \qe@origexamstatement{#1}%
  }%
}
\makeatother


\title{\textbf{QUALIFYING EXAM: APPLIED MATHEMATICS}\\[0.5em]\Large SPRING 2025}
\date{}
\author{}

\begin{document}

\maketitle
\vspace{-3em}

\tableofcontents

\section{Statement of the problems}

\examstatement{Problem 1. [15 points]}
Assume that \(f\) is a periodic function with period \(2\pi\). Considering the integration \(I(f) = \int_{0}^{2\pi} f(x)\,\mathrm{d}x\), prove the following statements on the order of convergence of composite trapezoidal rule: \(I_{h}^{Tr}(f) = h \sum_{i=1}^{n} f(i \cdot h)\), where \(h = 2\pi/n\).
\begin{enumerate}[label=(\roman*)]
    \item Assume further that \(f \in C^{m}([0, 2\pi])\), prove that the above rule achieves \(m\)-th order convergence, i.e.\ \(\exists C > 0\) s.t.\ \(\lvert I(f) - I_{h}^{Tr}(f)\rvert \le C \cdot h^{m}\).
    \item Assume further that \(f\) is analytic, prove that the above rule achieves exponential convergence.
\end{enumerate}

\examstatement{Problem 2. [15 points]}
Let \(f\) be a continuous function that is \(C^{m}\) (\(m \ge 1\)), such that \(f(\alpha) = f^{(1)}(\alpha) = \cdots = f^{(m-1)}(\alpha) = 0\) and \(f^{(m)} \neq 0\). Answer the following questions:
\begin{enumerate}[label=(\roman*)]
    \item If \(m = 1\), i.e.\ \(\alpha\) is a simple zero of the function \(f\), prove that there exists a neighborhood of \(\alpha\): \(D(\alpha, \delta) = \{x : |x - \alpha| \le \delta\}\), s.t.\ \(\forall x_{0} \in D(\alpha, \delta)\), the sequence generated by Newton's iteration converges to \(\alpha\), and the rate of convergence is quadratic.
    \item If \(m > 1\), i.e.\ \(\alpha\) is a zero with multiplicity greater than \(1\), what is the rate of convergence of Newton's iteration? How about the modified Newton's iteration:
    \[
    x_{k+1} = x_{k} - m \frac{f(x_{k})}{f'(x_{k})}?
    \]
    Please prove your results.
    \item If the multiplicity \(m\) is greater than \(1\) in general but the value is unknown, can you propose an iteration method that achieves superlinear convergence? (You don't need to prove the result.)
\end{enumerate}

\examstatement{Problem 3. [20 points]}
The Sherman--Morrison formula (or the generalized version named Sherman--Morrison--Woodbury formula) is an useful tool in numerical linear algebra. It states: if \(A \in \mathbb{R}^{n \times n}\) is an invertible matrix, \(u, v \in \mathbb{R}^{n}\) are two vectors satisfying \(1 + v^{T} A^{-1} u \neq 0\), then the matrix \(A + uv^{T}\) is invertible and can be computed by

\[
(A + uv^{T})^{-1} = A^{-1} - \frac{A^{-1} uv^{T} A^{-1}}{1 + v^{T} A^{-1} u}.
\]

\begin{enumerate}[label=(\roman*)]
    \item Prove the above statement.
    \item If \(\{s_{j}\}_{j=1}^{n}\) and \(\{t_{j}\}_{j=1}^{n}\) are two sets of points in \([0, 1]\), \(A \in \mathbb{R}^{n \times n}\) is a matrix defined by:
    \[
    A_{j,k} = 4n * \delta_{j,k} + \cos(t_{j} - s_{k}),
    \]
    where \(\delta_{j,k}\) is the Kronecker delta (\(\delta_{j,k} = 1\) when \(j = k\), \(\delta_{j,k} = 0\) otherwise). \(b \in \mathbb{R}^{n}\) is an arbitrary vector. Please give methods for the evaluation of \(Ab\), \(A^{-1}\), and \(|A|\), all with the computational complexity of \(O(n)\).
\end{enumerate}

\examstatement{Problem 4. [15 points]}
For the system

\[
u_{t} = -v_{xx}, \qquad v_{t} = u_{xx},
\]

analyse the truncation error and stability of the scheme

\[
\frac{u_{j}^{n+1} - u_{j}^{n}}{\tau} = -\frac{1}{2h^{2}}\bigl(v_{j+1}^{n} - 2v_{j}^{n} + v_{j-1}^{n} + v_{j+1}^{n+1} - 2v_{j}^{n+1} + v_{j-1}^{n+1}\bigr),
\]

\[
\frac{v_{j}^{n+1} - v_{j}^{n}}{\tau} = \frac{1}{2h^{2}}\bigl(u_{j+1}^{n} - 2u_{j}^{n} + u_{j-1}^{n} + u_{j+1}^{n+1} - 2u_{j}^{n+1} + u_{j-1}^{n+1}\bigr).
\]

\examstatement{Problem 5. [15 points]}
For the advection equation \(u_{t} + a u_{x} = 0\) (\(a\) is a constant), the Lax--Wendroff scheme reads as

\[
u_{j}^{n+1} = -\frac12 \nu(1-\nu) u_{j+1}^{n} + (1-\nu^{2}) u_{j}^{n} + \frac12 \nu(1+\nu) u_{j-1}^{n},
\]

where \(\nu = a\tau/h\).
\begin{enumerate}[label=(\roman*)]
    \item For the Cauchy problem imposed on the real line, show that
    \[
    \|u^{n+1}\|_{2}^{2} = \|u^{n}\|_{2}^{2} - \frac12 \nu^{2}(1-\nu^{2}) \bigl( \|\delta_{x}^{+} u^{n}\|_{2}^{2} - \langle \delta_{x}^{+} u^{n}, \delta_{x}^{-} u^{n} \rangle \bigr),
    \]
    where
    \[
    \|v\|_{2}^{2} = \sum_{j} |v_{j}|^{2}, \quad \langle v, w \rangle = \sum_{j} v_{j} w_{j}, \quad \delta_{x}^{-} v_{j} = v_{j} - v_{j-1}, \quad \delta_{x}^{+} v_{j} = v_{j+1} - v_{j}.
    \]
    \item Suppose \(a > 0\), for the problem imposed on \((0, 1)\) with homogeneous boundary condition at \(x = 0\) (i.e., \(u_{0}^{n} = 0\)), give a simple numerical boundary condition for \(x = 1\) such that the Lax--Wendroff scheme is stable.
\end{enumerate}

\noindent\textbf{Remark: You can choose either 6 or 7. The points will be decided as \(\max(6,7)\).}

\examstatement{Problem 6. [20 points]}
Consider a harmonic oscillator with a cubic damping term

\[
y'' + y + \epsilon (y')^{3} = 0, \qquad y(0) = 1,\quad y'(0) = 0,
\]

where \(y = y(t)\), \(t \ge 0\), \(\epsilon > 0\).
\begin{enumerate}[label=(\roman*)]
    \item For small \(\epsilon\), use the multiple-scale method to study the behavior of \(y(t)\) for large \(t\), i.e., construct a proper asymptotic solution.
    \item Make a conclusion on the validity of your asymptotic solution. Briefly justify your conclusion.
\end{enumerate}

\examstatement{Problem 7. [20 points]}
Let \(K\) be a nonempty closed convex set in \(\mathbb{R}^{n}\). For any \(z \in \mathbb{R}^{n}\), define \(\Pi_{K}(z)\) as the optimal solution of the problem \(\min \tfrac12 \|y - z\|\), s.t.\ \(y \in K\). Prove the following statements.
\begin{enumerate}[label=(\roman*)]
    \item Show that \(\Pi_{K}(z)\) satisfies
    \[
    \langle z - \Pi_{K}(z), d - \Pi_{K}(z) \rangle \le 0, \qquad \forall d \in K.
    \]
    \item Define
    \[
    \Theta(z) := \tfrac12 \|z - \Pi_{K}(z)\|^{2}.
    \]
    Prove that \(\Theta(\cdot)\) is continuous differentiable convex function and
    \[
    \nabla \Theta(z) = z - \Pi_{K}(z).
    \]
\end{enumerate}

\noindent\rule{\linewidth}{0.4pt}

\section{Statement and solutions}

\subsection{Problem 1. [15 points]}
Assume that \(f\) is a periodic function with period \(2\pi\). Considering the integration \(I(f) = \int_{0}^{2\pi} f(x)\,\mathrm{d}x\), prove the following statements on the order of convergence of composite trapezoidal rule: \(I_{h}^{Tr}(f) = h \sum_{i=1}^{n} f(i \cdot h)\), where \(h = 2\pi/n\).
\begin{enumerate}[label=(\roman*)]
    \item Assume further that \(f \in C^{m}([0, 2\pi])\), prove that the above rule achieves \(m\)-th order convergence, i.e.\ \(\exists C > 0\) s.t.\ \(\lvert I(f) - I_{h}^{Tr}(f)\rvert \le C \cdot h^{m}\).
    \item Assume further that \(f\) is analytic, prove that the above rule achieves exponential convergence.
\end{enumerate}

\setcounter{section}{1}
\setcounter{theorem}{0}
\setcounter{definition}{0}
\setcounter{remark}{0}
\setcounter{note}{0}

\subsubsection{Preliminaries: Euler--Maclaurin, after Qingshu}

The official CAM syllabus names this circle as interpolation and approximation, including trigonometric interpolation \cite{conte00}. Problem~1 is the periodic trapezoidal rule. Part~(i) is the statement that periodicity upgrades the usual \(O(h^{2})\) error to the full smoothness order \(O(h^{m})\). Part~(ii) is what that becomes when \(f\) is analytic: the error is exponential in \(n=2\pi/h\).

The piecewise-linear interpolation remainder does \emph{not} give (i). On each panel \([ih,(i+1)h]\) one has
\[
f(x)-p_{1}(x)
=
\frac{f''(\xi_{x})}{2}(x-ih)\bigl(x-(i+1)h\bigr),
\]
and integrating produces the classical composite estimate \(I(f)-I_{h}^{Tr}(f)=-\frac{2\pi}{12}h^{2}f''(\xi)=O(h^{2})\). Higher smoothness only improves the constant. The missing mechanism is cancellation of the endpoint corrections, which is exactly Euler--Maclaurin. The background below is Qingshu, \S8.5 (printed pp.~94--97; PDF pp.~102--107) \cite{qingshu11}.

\paragraph{Qiuzhen's quadrature questions, 2025--2026.}

{\footnotesize
\begin{center}
\begin{tabular}{@{}>{\raggedright\arraybackslash}p{2.45cm}
                   >{\raggedright\arraybackslash}p{4.35cm}
                   >{\raggedright\arraybackslash}p{3.45cm}
                   >{\raggedright\arraybackslash}p{3.55cm}@{}}
\toprule
Paper & Typical ask & Machine & Relation to 2025 Spring \\
\midrule
2025 Spring P1 & periodic trapezoid, order \(m\) / exponential & Euler--Maclaurin; Fourier aliasing & this problem \\
2025 Fall P2 & Gauss--Chebyshev nodes and weights & interpolatory quadrature at zeros of \(T_{n+1}\) & same word ``quadrature'', different nodes \\
2026 Spring P2 & Gauss--Legendre error via \(\omega^{2}\) & Hermite remainder / Peano & Qingshu Ch.~10, not Ch.~8 \\
\bottomrule
\end{tabular}
\end{center}
}

\begin{definition}[Periodic \(C^{m}\), and the two trapezoidal sums]
\label{def:p1-periodic}
A \(2\pi\)-periodic function of class \(C^{m}\) is a function \(f\in C^{m}(\mathbb{R})\) with \(f(x+2\pi)=f(x)\). Equivalently, \(f\in C^{m}([0,2\pi])\) and
\[
f^{(j)}(0)
=
f^{(j)}(2\pi),
\qquad
j=0,1,\dots,m.
\]
(The exam's phrase ``periodic'' \(+\) \(f\in C^{m}([0,2\pi])\) is read this way: without matching jets at the endpoints the periodic extension is not \(C^{m}\), and the order in (i) drops to the first unmatched derivative.) Write
\[
I(f)
\coloneqq
\int_{0}^{2\pi} f(x)\,\mathrm{d}x,
\qquad
h
=
\frac{2\pi}{n},
\qquad
I_{h}^{Tr}(f)
\coloneqq
h\sum_{i=1}^{n} f(ih).
\]
The ordinary composite trapezoidal sum on \([0,2\pi]\) is
\begin{equation}
\label{eq:p1-ordinary-trap}
T_{h}(f)
\coloneqq
\frac{h}{2}\bigl(f(0)+f(2\pi)\bigr)
+
h\sum_{i=1}^{n-1} f(ih).
\end{equation}
If \(f(0)=f(2\pi)\), then \(I_{h}^{Tr}(f)=T_{h}(f)\). The exam writes the periodic form; the Euler--Maclaurin identity is most cleanly stated for \(T_{h}\).
\end{definition}

\paragraph{Qingshu's periodic Bernoulli functions.}

Qingshu does not start from the classical Bernoulli numbers. The functions \(b_{n}\) are defined by integrating the sawtooth and imposing mean zero; the values \(b_{n}(0)\) then become the coefficients in Euler--Maclaurin.

\begin{definition}[Periodic functions \(b_{n}\); Qingshu, Definition~8.1]
\label{def:p1-bn}
Let \(b_{1}\) be the \(1\)-periodic extension of \(x-\tfrac12\) from \([0,1)\). For each \(n\ge 2\), let \(b_{n}\) be the unique \(1\)-periodic antiderivative of \(b_{n-1}\) with
\[
\int_{0}^{1} b_{n}(x)\,\mathrm{d}x
=
0.
\]
Thus \(b_{n}'=b_{n-1}\) almost everywhere, and \(b_{n}(x)=B_{n}(\{x\})/n!\) in the usual Bernoulli-polynomial notation. On \([0,1)\) the first four are
\begin{align*}
b_{1}(x)
&=
x-\tfrac12,
&
b_{2}(x)
&=
\tfrac12 x^{2}-\tfrac12 x+\tfrac{1}{12},
\\
b_{3}(x)
&=
\tfrac16 x^{3}-\tfrac14 x^{2}+\tfrac{1}{12}x,
&
b_{4}(x)
&=
\tfrac{1}{24}x^{4}-\tfrac{1}{12}x^{3}+\tfrac{1}{24}x^{2}-\tfrac{1}{720}.
\end{align*}
In particular \(b_{2}(0)=\frac{1}{12}\) and \(b_{4}(0)=-\frac{1}{720}\).
\end{definition}

\begin{proposition}[Fourier series of \(b_{n}\); Qingshu, Proposition~8.4]
\label{prop:p1-bn-fourier}
For \(k\in\mathbb{N}\),
\begin{align*}
b_{1}(x)
&\sim
-\frac{1}{\pi}\sum_{n=1}^{\infty}\frac{\sin(2n\pi x)}{n},
\\
b_{2k}(x)
&=
\frac{2(-1)^{k+1}}{(2\pi)^{2k}}\sum_{n=1}^{\infty}\frac{\cos(2n\pi x)}{n^{2k}},
\\
b_{2k+1}(x)
&=
\frac{2(-1)^{k+1}}{(2\pi)^{2k+1}}\sum_{n=1}^{\infty}\frac{\sin(2n\pi x)}{n^{2k+1}}.
\end{align*}
The series for \(b_{1}\) converges to \(b_{1}(x)\) off \(\mathbb{Z}\) and to \(0\) at integers (Dirichlet).
\end{proposition}

\begin{corollary}[Values and bounds; Qingshu, Corollary~8.2]
\label{cor:p1-bn-bound}
For every \(k\in\mathbb{N}\),
\[
b_{2k+1}(0)
=
0,
\qquad
\lvert b_{k}(x)\rvert
\le
\frac{\zeta(k)}{(2\pi)^{k}}\qquad (k\ge 2),
\]
and \(\lvert b_{1}(x)\rvert\le\tfrac12\).
\end{corollary}

\paragraph{Qingshu's Euler--Maclaurin formula.}

The unit-step identity is the one printed in the notes. The exam needs the same identity after the change of scale \(x=th\).

\begin{proposition}[Order-zero formula; Qingshu, Proposition~8.3 / (8.41)]
\label{prop:p1-em0}
Let \(a,b\in\mathbb{Z}\) and \(f\in C^{1}[a,b]\). Then
\begin{equation}
\label{eq:p1-em0}
\sum_{k=a}^{b} f(k)
=
\int_{a}^{b} f(x)\,\mathrm{d}x
+
\frac{f(a)+f(b)}{2}
+
\int_{a}^{b} b_{1}(x)\,f'(x)\,\mathrm{d}x.
\end{equation}
\end{proposition}

\begin{theorem}[Euler--Maclaurin; Qingshu, Theorem~8.8 / (8.42)]
\label{thm:p1-em}
Let \(a,b,m\in\mathbb{Z}\) with \(m\ge 2\), and let \(f\in C^{m}[a,b]\). Then
\begin{equation}
\label{eq:p1-em}
\sum_{k=a}^{b} f(k)
-
\int_{a}^{b} f(x)\,\mathrm{d}x
=
\frac{f(a)+f(b)}{2}
+
\sum_{k=2}^{m}
\bigl[f^{(k-1)}(b)-f^{(k-1)}(a)\bigr]\,b_{k}(0)
+
(-1)^{m+1}
\int_{a}^{b} b_{m}(x)\,f^{(m)}(x)\,\mathrm{d}x.
\end{equation}
\end{theorem}

\begin{proof}[Proof, after Qingshu]
The case \(m=1\) is \eqref{eq:p1-em0}. For \(m=2\), integrate by parts using \(b_{2}'=b_{1}\) and the \(1\)-periodicity of \(b_{2}\) (so \(b_{2}(b)=b_{2}(a)=b_{2}(0)\)):
\[
\int_{a}^{b} b_{1}(x)\,f'(x)\,\mathrm{d}x
=
\bigl[f'(b)-f'(a)\bigr]\,b_{2}(0)
-
\int_{a}^{b} b_{2}(x)\,f''(x)\,\mathrm{d}x.
\]
Substitute into \eqref{eq:p1-em0}. The general formula is the same integration by parts, repeated.
\end{proof}

\begin{corollary}[Even form; Qingshu, Corollary~8.3 / (8.44)]
\label{cor:p1-em-even}
If \(f\in C^{2m}[a,b]\), then odd endpoint terms drop by \(b_{2k+1}(0)=0\), and
\begin{equation}
\label{eq:p1-em-even}
\sum_{k=a}^{b} f(k)
-
\int_{a}^{b} f(x)\,\mathrm{d}x
=
\frac{f(a)+f(b)}{2}
+
\sum_{k=1}^{m}
\bigl[f^{(2k-1)}(b)-f^{(2k-1)}(a)\bigr]\,b_{2k}(0)
-
\int_{a}^{b} b_{2m}(x)\,f^{(2m)}(x)\,\mathrm{d}x.
\end{equation}
\end{corollary}

\paragraph{How to use it on the exam interval.}

Qingshu sums \(f(k)\) with step \(1\). The exam samples \(f(ih)\) with step \(h\). The reduction is the substitution \(F(t)=f(th)\).

\begin{lemma}[Euler--Maclaurin at step \(h\)]
\label{lem:p1-em-h}
Let \(h=2\pi/n\) and \(f\in C^{m}([0,2\pi])\), with \(m\ge 1\). Then
\begin{equation}
\label{eq:p1-em-h}
T_{h}(f)-I(f)
=
\sum_{k=2}^{m}
h^{k}
\bigl[f^{(k-1)}(2\pi)-f^{(k-1)}(0)\bigr]\,b_{k}(0)
+
R_{m}(f,h),
\end{equation}
where the \(m=1\) sum is empty, and
\begin{equation}
\label{eq:p1-Rm}
R_{m}(f,h)
=
(-1)^{m+1}
h^{m}
\int_{0}^{2\pi}
b_{m}\!\left(\frac{x}{h}\right)
f^{(m)}(x)\,
\mathrm{d}x.
\end{equation}
In particular
\begin{equation}
\label{eq:p1-Rm-bound}
\bigl\lvert R_{m}(f,h)\bigr\rvert
\le
C_{m}\,h^{m}\,\lVert f^{(m)}\rVert_{L^{1}[0,2\pi]},
\qquad
C_{1}
=
\tfrac12,
\qquad
C_{m}
=
\frac{\zeta(m)}{(2\pi)^{m}}\quad (m\ge 2).
\end{equation}
\end{lemma}

\begin{proof}
Set \(F(t)\coloneqq f(th)\) for \(t\in[0,n]\), so \(F^{(j)}(t)=h^{j}f^{(j)}(th)\). Apply \cref{prop:p1-em0} if \(m=1\), and \cref{thm:p1-em} if \(m\ge 2\), on the integers from \(0\) to \(n\):
\begin{align*}
\sum_{k=0}^{n} F(k)
-
\int_{0}^{n} F(t)\,\mathrm{d}t
&=
\frac{F(0)+F(n)}{2}
+
\sum_{k=2}^{m}
\bigl[F^{(k-1)}(n)-F^{(k-1)}(0)\bigr]\,b_{k}(0)
\\
&\qquad
+
(-1)^{m+1}
\int_{0}^{n} b_{m}(t)\,F^{(m)}(t)\,\mathrm{d}t.
\end{align*}
The integral on the left is \(h^{-1}I(f)\). Multiply through by \(h\). The left-hand side becomes
\[
h\sum_{k=0}^{n} f(kh)
-
I(f)
=
T_{h}(f)
+
\frac{h}{2}\bigl(f(0)+f(2\pi)\bigr)
-
I(f),
\]
and the displayed endpoint sum on the right is exactly \(\frac{h}{2}(f(0)+f(2\pi))\) plus the Bernoulli corrections. What remains is \(T_{h}(f)-I(f)\) as in \eqref{eq:p1-em-h}. The remainder integral is
\[
(-1)^{m+1}h^{m+1}
\int_{0}^{n}
b_{m}(t)\,f^{(m)}(th)\,\mathrm{d}t.
\]
The substitution \(x=th\) produces \eqref{eq:p1-Rm}. The bound \eqref{eq:p1-Rm-bound} is \cref{cor:p1-bn-bound}.
\end{proof}

\begin{remark}[What periodicity does]
\label{rem:p1-cancel}
Every coefficient in the sum \eqref{eq:p1-em-h} is a jump \(f^{(k-1)}(2\pi)-f^{(k-1)}(0)\). On a non-periodic interval those jumps are the reason the trapezoidal rule stays \(O(h^{2})\) no matter how smooth \(f\) is (the first surviving term is \(h^{2}b_{2}(0)(f'(2\pi)-f'(0))\)). If \(f\) is \(2\pi\)-periodic of class \(C^{m}\), every jump through order \(m-1\) vanishes, and \eqref{eq:p1-em-h} collapses to \(T_{h}(f)-I(f)=R_{m}(f,h)=O(h^{m})\). That is the whole of part~(i).
\end{remark}

\subsubsection{Solution of (i)}

\begin{theorem}[\(C^{m}\) periodic \(\Rightarrow\) order \(m\)]
\label{thm:p1-order-m}
Let \(f\) be \(2\pi\)-periodic of class \(C^{m}\), \(m\ge 1\). Then \(I_{h}^{Tr}(f)=T_{h}(f)\) and
\[
\bigl\lvert I(f)-I_{h}^{Tr}(f)\bigr\rvert
\le
C_{m}\,h^{m}\,\lVert f^{(m)}\rVert_{L^{1}[0,2\pi]},
\]
with \(C_{m}\) as in \eqref{eq:p1-Rm-bound}. In particular the exam bound holds with \(C=C_{m}\lVert f^{(m)}\rVert_{L^{1}}\).
\end{theorem}

\begin{proof}
Periodicity gives \(f(0)=f(2\pi)\), hence \(I_{h}^{Tr}=T_{h}\) by \cref{def:p1-periodic}. The same periodicity gives \(f^{(j)}(0)=f^{(j)}(2\pi)\) for \(0\le j\le m\), so every Bernoulli endpoint term in \cref{lem:p1-em-h} vanishes. The identity reduces to \(I_{h}^{Tr}(f)-I(f)=R_{m}(f,h)\), and \eqref{eq:p1-Rm-bound} is the claim.
\end{proof}

\begin{remark}[Fourier aliasing, and why it is not the proof of (i) for \(m=1\)]
\label{rem:p1-alias-cm}
Write \(\hat f(k)=(2\pi)^{-1}\int_{0}^{2\pi}f(x)e^{-ikx}\,\mathrm{d}x\). If the Fourier series of \(f\) may be summed absolutely,
\begin{equation}
\label{eq:p1-alias}
I(f)
=
2\pi\,\hat f(0),
\qquad
I_{h}^{Tr}(f)
=
2\pi\sum_{\ell\in\mathbb{Z}}\hat f(\ell n),
\end{equation}
hence
\[
I(f)-I_{h}^{Tr}(f)
=
-2\pi\sum_{\ell\neq 0}\hat f(\ell n).
\]
(The inner geometric sum \(h\sum_{j=1}^{n}e^{ikjh}\) equals \(2\pi\) if \(n\mid k\) and equals \(0\) otherwise.) Integration by parts on the circle gives \(\lvert\hat f(k)\rvert\le (2\pi)^{-1}\lvert k\rvert^{-m}\lVert f^{(m)}\rVert_{L^{1}}\). For \(m\ge 2\) one may sum the aliases:
\[
\bigl\lvert I(f)-I_{h}^{Tr}(f)\bigr\rvert
\le
n^{-m}\,\lVert f^{(m)}\rVert_{L^{1}}\cdot 2\zeta(m)
=
O(h^{m}).
\]
For \(m=1\) the comparison series \(\sum_{\ell\neq 0}\lvert\ell\rvert^{-1}\) diverges, so this bound does not close. Euler--Maclaurin is the argument that covers every \(m\ge 1\), including the exam's \(C^{1}\) case.
\end{remark}

\subsubsection{Solution of (ii)}

``Analytic'' for a \(2\pi\)-periodic function means the periodic extension is real-analytic on \(\mathbb{R}\). Equivalently, \(f\) extends holomorphically to a strip \(\lvert\operatorname{Im}z\rvert<\alpha\) for some \(\alpha>0\).

\begin{lemma}[Exponential decay of Fourier coefficients]
\label{lem:p1-fourier-exp}
If \(f\) is \(2\pi\)-periodic and holomorphic in \(\lvert\operatorname{Im}z\rvert<\alpha\), then for every \(\beta\in(0,\alpha)\) there is \(M_{\beta}<\infty\) such that
\[
\lvert\hat f(k)\rvert
\le
M_{\beta}\,e^{-\beta\lvert k\rvert},
\qquad
k\in\mathbb{Z}.
\]
\end{lemma}

\begin{proof}
Shift the integration contour of \(\hat f(k)\) to \(\operatorname{Im}z=\beta\operatorname{sign}(-k)\) if \(k\neq 0\) (justified by periodicity in the real direction and holomorphy in the strip). Then \(\lvert e^{-ikz}\rvert=e^{-\beta\lvert k\rvert}\), and \(M_{\beta}\) is \((2\pi)^{-1}\) times the maximum of \(\lvert f\rvert\) on the two lines \(\operatorname{Im}z=\pm\beta\), times the length \(2\pi\).
\end{proof}

\begin{theorem}[Analytic periodic \(\Rightarrow\) exponential convergence]
\label{thm:p1-exp}
Let \(f\) be \(2\pi\)-periodic and holomorphic in a strip of width \(\alpha>0\). Then for every \(\beta\in(0,\alpha)\) there is \(C_{\beta}<\infty\) such that
\[
\bigl\lvert I(f)-I_{h}^{Tr}(f)\bigr\rvert
\le
C_{\beta}\,e^{-\beta n}
=
C_{\beta}\,e^{-2\pi\beta/h}.
\]
\end{theorem}

\begin{proof}
Absolute convergence of the Fourier series is \cref{lem:p1-fourier-exp}, so the aliasing identity \eqref{eq:p1-alias} applies and
\[
\bigl\lvert I(f)-I_{h}^{Tr}(f)\bigr\rvert
\le
2\pi\sum_{\ell\neq 0}
M_{\beta}\,e^{-\beta n\lvert\ell\rvert}
=
\frac{4\pi M_{\beta}\,e^{-\beta n}}{1-e^{-\beta n}}.
\]
The right-hand side is \(O(e^{-\beta n})\).
\end{proof}

\begin{remark}[The same conclusion from Euler--Maclaurin alone]
\label{rem:p1-em-exp}
If one prefers not to name Fourier coefficients, apply \cref{thm:p1-order-m} at every order. Cauchy's estimate on a circle of radius \(r\in(0,\alpha)\) in the complex plane gives \(\lVert f^{(m)}\rVert_{\infty}\le M\,m!\,r^{-m}\). Combined with \eqref{eq:p1-Rm-bound},
\[
\bigl\lvert I(f)-I_{h}^{Tr}(f)\bigr\rvert
\le
C\,\Bigl(\frac{h}{2\pi r}\Bigr)^{m}m!.
\]
Stirling's bound \(m!\le (m/e)^{m}\sqrt{2\pi m}\,e^{1/(12m)}\) and the choice \(m\sim 2\pi r/h\) produce an estimate \(O(e^{-c/h})\) for some \(c>0\). This is the same exponential, read off Qingshu's remainder rather than off aliasing. The Fourier argument is shorter.
\end{remark}

\begin{examnote}[What the two parts are]
\label{examnote:p1}
\begin{enumerate}[label=\textup{(\roman*)}]
    \item Quote Qingshu's formula at step \(h\), \cref{lem:p1-em-h}. Periodicity kills every jump \(f^{(k-1)}(2\pi)-f^{(k-1)}(0)\). The error is the integral remainder \(R_{m}=O(h^{m})\). Do not stop at the piecewise-linear interpolation remainder: that is only \(O(h^{2})\).
    \item Analytic \(\Rightarrow\) holomorphic in a strip \(\Rightarrow\) \(\lvert\hat f(k)\rvert\le M e^{-\beta\lvert k\rvert}\). The trapezoidal sum aliases to \(2\pi\sum_{\ell}\hat f(\ell n)\), so the error is \(O(e^{-\beta n})\). The Euler--Maclaurin-only route is \cref{rem:p1-em-exp}.
\end{enumerate}
\end{examnote}

\begin{remark}[Local Qingshu slice]
\label{rem:p1-qingshu-em}
The statements copied above are Qingshu, \S8.5.1, Theorem~8.8 and Corollary~8.3 \cite{qingshu11}. A six-page excerpt (original PDF pp.~102--107) sits at \texttt{qe-review/applied-math/reference/qingshu-euler-maclaurin/}.
\end{remark}

% Restore parent section so subsequent \subsection headings stay under
% ``Statement and solutions''.
\setcounter{section}{2}


\subsection{Problem 2. [15 points]}
Let \(f\) be a continuous function that is \(C^{m}\) (\(m \ge 1\)), such that \(f(\alpha) = f^{(1)}(\alpha) = \cdots = f^{(m-1)}(\alpha) = 0\) and \(f^{(m)} \neq 0\). Answer the following questions:
\begin{enumerate}[label=(\roman*)]
    \item If \(m = 1\), i.e.\ \(\alpha\) is a simple zero of the function \(f\), prove that there exists a neighborhood of \(\alpha\): \(D(\alpha, \delta) = \{x : |x - \alpha| \le \delta\}\), s.t.\ \(\forall x_{0} \in D(\alpha, \delta)\), the sequence generated by Newton's iteration converges to \(\alpha\), and the rate of convergence is quadratic.
    \item If \(m > 1\), i.e.\ \(\alpha\) is a zero with multiplicity greater than \(1\), what is the rate of convergence of Newton's iteration? How about the modified Newton's iteration:
    \[
    x_{k+1} = x_{k} - m \frac{f(x_{k})}{f'(x_{k})}?
    \]
    Please prove your results.
    \item If the multiplicity \(m\) is greater than \(1\) in general but the value is unknown, can you propose an iteration method that achieves superlinear convergence? (You don't need to prove the result.)
\end{enumerate}

\setcounter{section}{2}
\setcounter{theorem}{0}
\setcounter{definition}{0}
\setcounter{remark}{0}
\setcounter{note}{0}

\subsubsection{Preliminaries: Newton's iteration, and the theorems the three parts invoke}

The official CAM syllabus names this circle as iterative methods for nonlinear equations: bisection, Newton, quasi-Newton, and fixed-point iteration, together with their convergence analysis \cite{conte00}. Problem~2 is the Newton package in one question. Part~(i) is the local quadratic theorem at a simple zero (the same statement as 2023 Fall Problem~1, without the secant companion). Part~(ii) is what that theorem becomes when \(f'(\alpha)=0\). Part~(iii) is the standard repair when the multiplicity is unknown.

The iteration is not defined in the exam statement. The three parts are unintelligible until it is, and until ``rate of convergence'' is a precise limit rather than a synonym of ``fast''.

\paragraph{Qiuzhen's Newton questions, 2023--2026.}

{\footnotesize
\begin{center}
\begin{tabular}{@{}>{\raggedright\arraybackslash}p{2.45cm}
                   >{\raggedright\arraybackslash}p{4.35cm}
                   >{\raggedright\arraybackslash}p{3.45cm}
                   >{\raggedright\arraybackslash}p{3.55cm}@{}}
\toprule
Paper & Typical ask & Machine & Relation to 2025 \\
\midrule
2023 Fall P1 & simple zero: quadratic limit; secant: product of two errors & Taylor of \(f\) at \(\alpha\); \(g'(\alpha)=0\) & (i) is this theorem \\
2024 Fall P1 & Newton on a concrete transcendental, with a prescribed start & same local theorem, plus a basin argument & not needed here \\
2025 Spring P2 & simple / multiple / unknown multiplicity & factorization \(f=(x-\alpha)^m h\) & this problem \\
2026 Spring P1 & derivative-free one-step map, order \(2\) & same \(Q\)-order machine, different \(\varphi\) & Newton is the model case \\
\bottomrule
\end{tabular}
\end{center}
}

\begin{definition}[Newton's iteration]
\label{def:newton}
Let \(f\) be differentiable on an interval, and suppose \(f'(x)\neq 0\) at the points that arise. Newton's iteration for a root of \(f(x)=0\), started at \(x_0\), is the recurrence
\begin{equation}
\label{eq:newton}
x_{k+1}
=
x_k
-
\frac{f(x_k)}{f'(x_k)},
\qquad
k=0,1,2,\dots.
\end{equation}
Equivalently, \(x_{k+1}\) is the unique root of the tangent line to \(y=f(x)\) at \(x_k\):
\[
y
=
f(x_k)
+
f'(x_k)\,(x-x_k).
\]
Setting \(y=0\) and solving for \(x\) recovers \eqref{eq:newton}. The map is well-defined at \(x_k\) if and only if \(f'(x_k)\neq 0\).
\end{definition}

The geometric content of \cref{def:newton} is that one replaces \(f\) by its first-order Taylor polynomial at the current guess and solves that linear equation exactly. If the graph is locally well approximated by its tangent, the next intercept is much closer to the root than the current guess; that is the source of quadratic convergence at a simple zero, and it is also why a flat tangent (\(f'(\alpha)=0\)) destroys the rate.

\begin{center}
\begin{tikzpicture}[>=Latex, scale=1.05]
  \draw[->] (-0.25,0) -- (5.15,0) node[right] {$x$};
  \draw[->] (0,-1.55) -- (0,2.65) node[above] {$y$};
  \draw[thick, domain=0.35:4.55, samples=120]
    plot (\x, {0.22*(\x-0.3)*(\x-0.3)-1.35});
  \node[font=\small] at (4.42,2.05) {$y=f(x)$};
  % simple root of 0.22(x-0.3)^2-1.35=0
  \fill (2.777,0) circle (1.5pt);
  \node[below] at (2.777,-0.08) {\(\alpha\)};
  \pgfmathsetmacro{\xzero}{4.20}
  \pgfmathsetmacro{\yzero}{0.22*(\xzero-0.3)*(\xzero-0.3)-1.35}
  \pgfmathsetmacro{\fpzero}{0.44*(\xzero-0.3)}
  \pgfmathsetmacro{\xone}{\xzero-\yzero/\fpzero}
  \fill (\xzero,0) circle (1.5pt);
  \node[below] at (\xzero,-0.08) {\(x_0\)};
  \draw[densely dashed] (\xzero,0) -- (\xzero,\yzero);
  \fill (\xzero,\yzero) circle (1.5pt);
  \draw[thick] (2.45, {\yzero+\fpzero*(2.45-\xzero)}) -- (4.50, {\yzero+\fpzero*(4.50-\xzero)});
  \fill (\xone,0) circle (1.5pt);
  \node[below] at (\xone,-0.08) {\(x_1\)};
\end{tikzpicture}
\end{center}

\begin{notation}[Error, and the Newton map]
\label{not:newton-map}
Write \(e_k\coloneqq x_k-\alpha\). Whenever \(f'(x)\neq 0\), the iteration is the fixed-point map
\[
g(x)
\coloneqq
x
-
\frac{f(x)}{f'(x)},
\qquad
x_{k+1}=g(x_k).
\]
A root \(\alpha\) of \(f\) is a fixed point of \(g\) as soon as \(g\) extends continuously through \(\alpha\). At a simple root this is automatic: \(f'(\alpha)\neq 0\), so \(g(\alpha)=\alpha\). At a multiple root \(f'(\alpha)=0=f(\alpha)\), and \(g(\alpha)\) is an indeterminate form; the continuous extension is computed from the factorization in \cref{def:multiplicity} and equals \(\alpha\) still.
\end{notation}

\begin{definition}[Multiplicity]
\label{def:multiplicity}
Let \(f\) be \(C^m\) near \(\alpha\), with \(m\ge 1\). The point \(\alpha\) is a root of multiplicity \(m\) if
\[
f(\alpha)
=
f'(\alpha)
=
\cdots
=
f^{(m-1)}(\alpha)
=
0
\qquad\text{and}\qquad
f^{(m)}(\alpha)\neq 0.
\]
Equivalently, one may factor
\[
f(x)
=
(x-\alpha)^m h(x)
\]
in a neighborhood of \(\alpha\), with \(h\) continuous and \(h(\alpha)\neq 0\). (If \(f\in C^{m+1}\) then \(h\in C^1\) and \(h(\alpha)=f^{(m)}(\alpha)/m!\).) The exam's standing hypothesis is exactly this factorization. The root is \emph{simple} when \(m=1\), i.e.\ when \(f(\alpha)=0\) and \(f'(\alpha)\neq 0\).
\end{definition}

\begin{definition}[\(Q\)-order of convergence]
\label{def:q-order}
Let \(x_k\to\alpha\) with \(x_k\neq\alpha\) for all \(k\). The convergence is
\begin{enumerate}[label=\textup{(\roman*)}]
    \item \emph{\(Q\)-linear} with rate \(\mu\in(0,1)\) if
    \[
    \lim_{k\to\infty}
    \frac{\lvert x_{k+1}-\alpha\rvert}{\lvert x_k-\alpha\rvert}
    =
    \mu;
    \]
    \item \emph{\(Q\)-superlinear} if the same limit exists and equals \(0\);
    \item of \emph{\(Q\)-order \(p>1\)} if
    \[
    \lim_{k\to\infty}
    \frac{x_{k+1}-\alpha}{(x_k-\alpha)^p}
    =
    C
    \in\mathbb{R}\setminus\{0\}.
    \]
\end{enumerate}
The case \(p=2\) is \emph{quadratic}. A weaker, and often more convenient, form of order \(p\) is the existence of \(M,\delta>0\) such that \(\lvert x-\alpha\rvert<\delta\) implies \(\lvert g(x)-\alpha\rvert\le M\lvert x-\alpha\rvert^p\); combined with the restriction \(M\delta^{p-1}\le\tfrac12\), this already forces local attraction and is what one proves first. The limit \(C\), when it exists, is the \emph{asymptotic error constant}.
\end{definition}

Linear convergence with rate \(\mu\) close to \(1\) is slow: errors shrink by a fixed fraction each step. Superlinear means that fraction itself tends to \(0\). Quadratic means the number of correct digits roughly doubles at each step, once one is already close. The exam's ``rate of convergence is quadratic'' in (i) is \cref{def:q-order}(iii) with \(p=2\), not merely \(x_k\to\alpha\).

The one-step theory is Ostrowski's theorem: the derivatives of \(g\) at the fixed point determine the order. Newton's method is the special case in which those derivatives can be read off from \(f\).

\begin{theorem}[Ostrowski]
\label{thm:ostrowski}
Let \(g\) be \(C^p\) near a fixed point \(\alpha\), with \(p\ge 1\). Write \(x_{k+1}=g(x_k)\).
\begin{enumerate}[label=\textup{(\roman*)}]
    \item If \(\lvert g'(\alpha)\rvert<1\), there is a neighborhood \(D(\alpha,\delta)=\{x:\lvert x-\alpha\rvert\le\delta\}\) such that \(x_0\in D(\alpha,\delta)\) implies \(x_k\to\alpha\). If \(g'(\alpha)\neq 0\), the convergence is \(Q\)-linear with rate \(\lvert g'(\alpha)\rvert\).
    \item If \(g'(\alpha)=\cdots=g^{(p-1)}(\alpha)=0\) and \(g^{(p)}(\alpha)\neq 0\), the convergence is of \(Q\)-order \(p\), with
    \[
    \lim_{k\to\infty}
    \frac{x_{k+1}-\alpha}{(x_k-\alpha)^p}
    =
    \frac{g^{(p)}(\alpha)}{p!},
    \]
    provided \(x_0\) is sufficiently close to \(\alpha\) and \(x_k\neq\alpha\).
\end{enumerate}
\end{theorem}

\begin{proof}
Taylor's formula with Lagrange remainder gives
\[
g(x)-\alpha
=
g(\alpha)
+
g'(\alpha)(x-\alpha)
+
\cdots
+
\frac{g^{(p-1)}(\alpha)}{(p-1)!}(x-\alpha)^{p-1}
+
\frac{g^{(p)}(\xi_x)}{p!}(x-\alpha)^p
\]
for some \(\xi_x\) between \(x\) and \(\alpha\). Under (i) one has \(g(x)-\alpha=g'(\alpha)(x-\alpha)+o(x-\alpha)\). Choose \(\delta\) so that \(\lvert x-\alpha\rvert\le\delta\) implies \(\lvert g(x)-\alpha\rvert\le q\lvert x-\alpha\rvert\) with \(\lvert g'(\alpha)\rvert<q<1\). Then \(D(\alpha,\delta)\) is invariant and \(\lvert e_k\rvert\le q^k\lvert e_0\rvert\to 0\). Dividing by \(e_k\) and passing to the limit yields the rate \(\lvert g'(\alpha)\rvert\) as soon as \(e_k\neq 0\) and \(g'\) is continuous.

Under (ii) the same expansion collapses to
\[
g(x)-\alpha
=
\frac{g^{(p)}(\xi_x)}{p!}(x-\alpha)^p.
\]
Continuity of \(g^{(p)}\) at \(\alpha\) produces a bound \(\lvert g(x)-\alpha\rvert\le M\lvert x-\alpha\rvert^p\) on a small ball. Shrinking the ball so that \(M\delta^{p-1}\le\tfrac12\) gives local attraction as in \cref{thm:newton-simple} below, and the displayed limit follows by sending \(\xi_{x_k}\to\alpha\).
\end{proof}

The hypothesis \(\lvert g'(\alpha)\rvert<1\) is sharp for attraction: if \(\lvert g'(\alpha)\rvert>1\) then \(\alpha\) is a repeller, and if \(\lvert g'(\alpha)\rvert=1\) the linearization does not decide. For Newton at a simple root one will find \(g'(\alpha)=0\), which is the strongest possible case of (i) and triggers (ii) with \(p=2\).

\begin{lemma}[Newton map at a simple root]
\label{lem:g-simple}
Suppose \(f\in C^2\) near \(\alpha\), with \(f(\alpha)=0\) and \(f'(\alpha)\neq 0\). Then \(g\) is \(C^1\) near \(\alpha\), \(g(\alpha)=\alpha\), and
\[
g'(\alpha)
=
0,
\qquad
g''(\alpha)
=
\frac{f''(\alpha)}{f'(\alpha)}.
\]
In particular Ostrowski's theorem applies with \(p=2\), and the asymptotic constant is \(f''(\alpha)/(2f'(\alpha))\).
\end{lemma}

\begin{proof}
Differentiating \(g(x)=x-f(x)/f'(x)\) on \(\{f'\neq 0\}\) gives
\[
g'(x)
=
\frac{f(x)f''(x)}{\bigl(f'(x)\bigr)^2}.
\]
At \(x=\alpha\) the numerator vanishes and the denominator does not, so \(g'(\alpha)=0\). A second differentiation, or L'H\^{o}pital on \(g'(x)\) as \(x\to\alpha\), yields \(g''(\alpha)=f''(\alpha)/f'(\alpha)\).
\end{proof}

\begin{theorem}[Newton: simple root, local quadratic convergence]
\label{thm:newton-simple}
Let \(f\) be \(C^2\) in a neighborhood of \(\alpha\), with \(f(\alpha)=0\) and \(f'(\alpha)\neq 0\). There exists \(\delta>0\) such that for every \(x_0\in D(\alpha,\delta)=\{x:\lvert x-\alpha\rvert\le\delta\}\), Newton's iteration \eqref{eq:newton} is well-defined, remains in \(D(\alpha,\delta)\), converges to \(\alpha\), and is quadratic: there is \(M>0\) with
\[
\bigl\lvert x_{k+1}-\alpha\bigr\rvert
\le
M\bigl\lvert x_k-\alpha\bigr\rvert^2.
\]
If \(x_k\neq\alpha\) for all \(k\), the asymptotic constant exists and equals
\begin{equation}
\label{eq:newton-const}
\lim_{k\to\infty}
\frac{x_{k+1}-\alpha}{(x_k-\alpha)^2}
=
\frac{f''(\alpha)}{2f'(\alpha)}.
\end{equation}
\end{theorem}

\begin{proof}
Since \(f'(\alpha)\neq 0\) and \(f'\) is continuous, there is \(\delta_0>0\) such that \(\lvert x-\alpha\rvert\le\delta_0\) implies \(f'(x)\neq 0\). Thus \(g\) is well-defined on \(D(\alpha,\delta_0)\).

Taylor-expand \(f\) at \(x_k\) and evaluate at \(\alpha\). For some \(\xi_k\) between \(x_k\) and \(\alpha\),
\[
0
=
f(\alpha)
=
f(x_k)
+
f'(x_k)(\alpha-x_k)
+
\tfrac12 f''(\xi_k)(\alpha-x_k)^2,
\]
hence
\[
\frac{f(x_k)}{f'(x_k)}
=
x_k-\alpha
-
\frac{f''(\xi_k)}{2f'(x_k)}(x_k-\alpha)^2.
\]
Subtracting from \(x_k\) produces the exact error formula
\begin{equation}
\label{eq:newton-lagrange}
x_{k+1}-\alpha
=
\frac{f''(\xi_k)}{2f'(x_k)}(x_k-\alpha)^2.
\end{equation}
On the compact ball \(D(\alpha,\delta_0)\) the continuous function \(\lvert f''\rvert/(2\lvert f'\rvert)\) attains a finite maximum \(M\). Therefore
\[
\bigl\lvert g(x)-\alpha\bigr\rvert
\le
M\bigl\lvert x-\alpha\bigr\rvert^2
\qquad\text{whenever}\qquad
\lvert x-\alpha\rvert\le\delta_0.
\]
Shrink to \(\delta\in(0,\delta_0]\) with \(M\delta\le\tfrac12\). If \(\lvert e_0\rvert\le\delta\), then \(\lvert e_1\rvert\le M\lvert e_0\rvert^2\le\tfrac12\lvert e_0\rvert\le\delta\), and by induction the orbit remains in \(D(\alpha,\delta)\), stays well-defined, and satisfies \(\lvert e_{k+1}\rvert\le\tfrac12\lvert e_k\rvert\). Hence \(e_k\to 0\). The same bound is the \(Q\)-order \(2\) estimate.

For the limit, send \(k\to\infty\) in \eqref{eq:newton-lagrange}: \(\xi_k\to\alpha\) and \(f'(x_k)\to f'(\alpha)\neq 0\), which is \eqref{eq:newton-const}. (If \(f''(\alpha)=0\) the order is at least \(2\), and may be higher; the exam's \(C^m\) hypothesis with \(m=1\) does not force \(f''(\alpha)\neq 0\), so ``quadratic'' is the guaranteed lower bound.)
\end{proof}

The neighborhood \(D(\alpha,\delta)\) is the object the exam asks one to produce. It is not unique, and no explicit \(\delta\) is required: existence follows from continuity of \(f'\) at a simple root, plus the smallness \(M\delta\le\tfrac12\) that turns a quadratic bound into a contraction of the error. The argument is local. Global convergence of Newton is a different theorem (Kantorovich), and is not asked.

\begin{lemma}[Newton map at a multiple root]
\label{lem:g-multiple}
Let \(f(x)=(x-\alpha)^m h(x)\) with \(m\ge 2\), \(h(\alpha)\neq 0\), and \(h\in C^1\) near \(\alpha\). Then, in a punctured neighborhood of \(\alpha\),
\[
\frac{f(x)}{f'(x)}
=
\frac{(x-\alpha)\,h(x)}{m h(x)+(x-\alpha)h'(x)},
\]
and the Newton map extends continuously through \(\alpha\) by \(g(\alpha)=\alpha\), with
\begin{equation}
\label{eq:gprime-mult}
g'(\alpha)
=
1-\frac{1}{m}.
\end{equation}
In particular \(\lvert g'(\alpha)\rvert=\lvert 1-1/m\rvert\in(0,1)\).
\end{lemma}

\begin{proof}
Differentiating the factorization gives
\[
f'(x)
=
(x-\alpha)^{m-1}\bigl(m h(x)+(x-\alpha)h'(x)\bigr).
\]
The displayed quotient for \(f/f'\) follows by cancelling \((x-\alpha)^{m-1}\), which is legal for \(x\neq\alpha\). The denominator at \(\alpha\) equals \(m h(\alpha)\neq 0\), so the right-hand side extends continuously, and
\[
g(x)-\alpha
=
(x-\alpha)
-
\frac{(x-\alpha)\,h(x)}{m h(x)+(x-\alpha)h'(x)}
=
(x-\alpha)\cdot
\frac{(m-1)h(x)+(x-\alpha)h'(x)}{m h(x)+(x-\alpha)h'(x)}.
\]
Dividing by \(x-\alpha\) and sending \(x\to\alpha\) yields \(g'(\alpha)=(m-1)/m=1-1/m\).
\end{proof}

\begin{theorem}[Newton: multiple root, \(Q\)-linear convergence]
\label{thm:newton-multiple}
Under the hypotheses of \cref{lem:g-multiple}, there is a neighborhood of \(\alpha\) in which Newton's iteration is well-defined and converges to \(\alpha\), and the convergence is \(Q\)-linear with rate
\[
\lim_{k\to\infty}
\frac{x_{k+1}-\alpha}{x_k-\alpha}
=
1-\frac{1}{m}.
\]
The rate is independent of \(h\), hence of \(f^{(m)}(\alpha)\). It tends to \(1\) as \(m\to\infty\): high multiplicity makes Newton arbitrarily slow.
\end{theorem}

\begin{proof}
\Cref{lem:g-multiple} supplies \(g'(\alpha)=1-1/m\in(0,1)\). Ostrowski's theorem~(\cref{thm:ostrowski}(i)) gives local attraction and the stated \(Q\)-linear rate. Well-definedness on a small ball follows from \(m h(x)+(x-\alpha)h'(x)\neq 0\) near \(\alpha\), which is the continuous denominator of \(f/f'\).
\end{proof}

The linearization \(e_{k+1}\approx\bigl(1-1/m\bigr)e_k\) is the whole of part~(ii) for unmodified Newton. For \(m=2\) the errors are (asymptotically) halved each step; for \(m=10\) they shrink by only \(10\%\). The quadratic picture of \cref{thm:newton-simple} is gone because the tangent at \(\alpha\) is horizontal, so \cref{def:newton} cannot even be evaluated at the root.

If the multiplicity is known, the standard repair is to take an \(m\)-fold Newton step.

\begin{definition}[Modified Newton iteration]
\label{def:modified-newton}
If \(\alpha\) is known to have multiplicity \(m\ge 1\), the modified Newton iteration is
\begin{equation}
\label{eq:modified-newton}
x_{k+1}
=
x_k
-
m\,\frac{f(x_k)}{f'(x_k)}.
\end{equation}
The associated map is \(g_m(x)=x-m f(x)/f'(x)\). For \(m=1\) this is ordinary Newton.
\end{definition}

\begin{theorem}[Modified Newton is quadratic]
\label{thm:newton-modified}
Under the hypotheses of \cref{lem:g-multiple}, the map \(g_m\) extends through \(\alpha\) by \(g_m(\alpha)=\alpha\), and
\[
g_m(x)-\alpha
=
\frac{(x-\alpha)^2 h'(x)}{m h(x)+(x-\alpha)h'(x)}.
\]
In particular \(g_m'(\alpha)=0\), and modified Newton is locally quadratic. If \(h'(\alpha)\neq 0\), the exact \(Q\)-order is \(2\), with asymptotic constant \(h'(\alpha)/(m h(\alpha))\).
\end{theorem}

\begin{proof}
From the quotient in \cref{lem:g-multiple},
\[
g_m(x)
=
x
-
m\cdot
\frac{(x-\alpha)\,h(x)}{m h(x)+(x-\alpha)h'(x)}.
\]
A direct rearrangement of the numerator, using \(x=\alpha+(x-\alpha)\), produces
\[
g_m(x)-\alpha
=
\frac{(x-\alpha)^2 h'(x)}{m h(x)+(x-\alpha)h'(x)},
\]
which is \(O\bigl((x-\alpha)^2\bigr)\) with a continuous coefficient. Thus \(g_m'(\alpha)=0\). If \(h'(\alpha)\neq 0\) the leading coefficient at \(\alpha\) is \(h'(\alpha)/(m h(\alpha))\neq 0\), and \cref{thm:ostrowski}(ii) with \(p=2\) applies. If \(h'(\alpha)=0\) the order is at least \(2\). Local attraction is the same \(M\delta\le\tfrac12\) argument as in \cref{thm:newton-simple}.
\end{proof}

The identity \(g_m'(\alpha)=0\) is the reason one multiplies by \(m\): ordinary Newton overshoots by the factor \(1-1/m\) still left in \(g'\); scaling the step by \(m\) cancels that factor exactly.

When \(m>1\) is unknown, \eqref{eq:modified-newton} cannot be formed. The standard device is to observe that the rational function \(u=f/f'\) has a \emph{simple} zero at \(\alpha\), independently of \(m\), and to apply ordinary Newton to \(u\).

\begin{proposition}[Newton on \(f/f'\), unknown multiplicity]
\label{prop:newton-u}
Let \(f(x)=(x-\alpha)^m h(x)\) with \(m\ge 1\), \(h(\alpha)\neq 0\), and \(h\in C^2\) near \(\alpha\). Set
\[
u(x)
\coloneqq
\frac{f(x)}{f'(x)}
=
\frac{(x-\alpha)\,h(x)}{m h(x)+(x-\alpha)h'(x)}
\qquad (x\neq\alpha),
\]
and extend by \(u(\alpha)=0\). Then \(u'(\alpha)=1/m\neq 0\), so \(\alpha\) is a simple zero of \(u\). Newton's iteration for \(u\),
\begin{equation}
\label{eq:schroeder}
x_{k+1}
=
x_k
-
\frac{u(x_k)}{u'(x_k)}
=
x_k
-
\frac{f(x_k)\,f'(x_k)}{\bigl(f'(x_k)\bigr)^2-f(x_k)f''(x_k)},
\end{equation}
is therefore locally quadratic at \(\alpha\) by \cref{thm:newton-simple}, hence in particular \(Q\)-superlinear. The second formula is the identity \(u'=1-f f''/(f')^2\).
\end{proposition}

\begin{proof}
The continuous extension \(u(\alpha)=0\) is \cref{lem:g-multiple}. Differentiating the extended quotient, or writing \(u(x)=(x-\alpha)/\bigl(m+(x-\alpha)h'(x)/h(x)\bigr)\) and evaluating the derivative at \(\alpha\), gives \(u'(\alpha)=1/m\neq 0\). The hypotheses of \cref{thm:newton-simple} hold for \(u\) in place of \(f\). Differentiating \(u=f/f'\) produces
\[
u'
=
\frac{(f')^2-f f''}{(f')^2},
\]
and substituting into \(x-u/u'\) is \eqref{eq:schroeder}.
\end{proof}

This is Schr\"oder's iteration (Newton applied to \(f/f'\)). It uses \(f''\), which the standing \(C^m\) hypothesis supplies for \(m\ge 2\). It is the natural answer to part~(iii): a single, explicitly written one-step map, with superlinear (in fact quadratic) local convergence, and with no a priori knowledge of \(m\).

\begin{examnote}[What the three parts are]
\label{examnote:p2}
\begin{enumerate}[label=\textup{(\roman*)}]
    \item \Cref{def:newton} plus \cref{thm:newton-simple}. Produce \(D(\alpha,\delta)\), check well-definedness from \(f'(\alpha)\neq 0\), and quote either the Lagrange formula \eqref{eq:newton-lagrange} or Ostrowski with \(g'(\alpha)=0\). The quadratic bound with \(M\delta\le\tfrac12\) is the attraction argument; the limit \eqref{eq:newton-const} is optional but matches the 2023 Fall wording.
    \item Unmodified Newton is \cref{thm:newton-multiple}: \(Q\)-linear, rate \(1-1/m\). Modified Newton \eqref{eq:modified-newton} is \cref{thm:newton-modified}: locally quadratic. Both proofs go through the factorization of \cref{def:multiplicity}, not through \eqref{eq:newton-lagrange} (the latter assumes \(f'(x_k)\neq 0\) at the root).
    \item Propose \eqref{eq:schroeder}. The one-line reason is that \(u=f/f'\) has a simple zero, so \cref{thm:newton-simple} applies to \(u\). The exam does not ask for the proof.
\end{enumerate}
\end{examnote}

\begin{remark}[What this is not]
\label{rem:newton-not}
The local theory above is not Kantorovich's theorem (a semiglobal existence statement from bounds on \(f''\) in a whole ball, producing a root that was not assumed). It is not a quasi-Newton method: the derivative is exact, not a secant or Broyden approximation (that is 2023 Fall P1(2) and 2026 Spring P1). It is not a statement about Newton's method for optimization, which is the same map applied to \(\nabla F=0\) and lives in the CAM optimization syllabus rather than in this question.
\end{remark}


\subsection{Problem 3. [20 points]}
The Sherman--Morrison formula (or the generalized version named Sherman--Morrison--Woodbury formula) is an useful tool in numerical linear algebra. It states: if \(A \in \mathbb{R}^{n \times n}\) is an invertible matrix, \(u, v \in \mathbb{R}^{n}\) are two vectors satisfying \(1 + v^{T} A^{-1} u \neq 0\), then the matrix \(A + uv^{T}\) is invertible and can be computed by

\[
(A + uv^{T})^{-1} = A^{-1} - \frac{A^{-1} uv^{T} A^{-1}}{1 + v^{T} A^{-1} u}.
\]

\begin{enumerate}[label=(\roman*)]
    \item Prove the above statement.
    \item If \(\{s_{j}\}_{j=1}^{n}\) and \(\{t_{j}\}_{j=1}^{n}\) are two sets of points in \([0, 1]\), \(A \in \mathbb{R}^{n \times n}\) is a matrix defined by:
    \[
    A_{j,k} = 4n * \delta_{j,k} + \cos(t_{j} - s_{k}),
    \]
    where \(\delta_{j,k}\) is the Kronecker delta (\(\delta_{j,k} = 1\) when \(j = k\), \(\delta_{j,k} = 0\) otherwise). \(b \in \mathbb{R}^{n}\) is an arbitrary vector. Please give methods for the evaluation of \(Ab\), \(A^{-1}\), and \(|A|\), all with the computational complexity of \(O(n)\).
\end{enumerate}

\subsection{Problem 4. [15 points]}
For the system

\[
u_{t} = -v_{xx}, \qquad v_{t} = u_{xx},
\]

analyse the truncation error and stability of the scheme

\[
\frac{u_{j}^{n+1} - u_{j}^{n}}{\tau} = -\frac{1}{2h^{2}}\bigl(v_{j+1}^{n} - 2v_{j}^{n} + v_{j-1}^{n} + v_{j+1}^{n+1} - 2v_{j}^{n+1} + v_{j-1}^{n+1}\bigr),
\]

\[
\frac{v_{j}^{n+1} - v_{j}^{n}}{\tau} = \frac{1}{2h^{2}}\bigl(u_{j+1}^{n} - 2u_{j}^{n} + u_{j-1}^{n} + u_{j+1}^{n+1} - 2u_{j}^{n+1} + u_{j-1}^{n+1}\bigr).
\]

\subsection{Problem 5. [15 points]}
For the advection equation \(u_{t} + a u_{x} = 0\) (\(a\) is a constant), the Lax--Wendroff scheme reads as

\[
u_{j}^{n+1} = -\frac12 \nu(1-\nu) u_{j+1}^{n} + (1-\nu^{2}) u_{j}^{n} + \frac12 \nu(1+\nu) u_{j-1}^{n},
\]

where \(\nu = a\tau/h\).
\begin{enumerate}[label=(\roman*)]
    \item For the Cauchy problem imposed on the real line, show that
    \[
    \|u^{n+1}\|_{2}^{2} = \|u^{n}\|_{2}^{2} - \frac12 \nu^{2}(1-\nu^{2}) \bigl( \|\delta_{x}^{+} u^{n}\|_{2}^{2} - \langle \delta_{x}^{+} u^{n}, \delta_{x}^{-} u^{n} \rangle \bigr),
    \]
    where
    \[
    \|v\|_{2}^{2} = \sum_{j} |v_{j}|^{2}, \quad \langle v, w \rangle = \sum_{j} v_{j} w_{j}, \quad \delta_{x}^{-} v_{j} = v_{j} - v_{j-1}, \quad \delta_{x}^{+} v_{j} = v_{j+1} - v_{j}.
    \]
    \item Suppose \(a > 0\), for the problem imposed on \((0, 1)\) with homogeneous boundary condition at \(x = 0\) (i.e., \(u_{0}^{n} = 0\)), give a simple numerical boundary condition for \(x = 1\) such that the Lax--Wendroff scheme is stable.
\end{enumerate}

\setcounter{section}{5}
\setcounter{theorem}{0}
\setcounter{definition}{0}
\setcounter{remark}{0}
\setcounter{note}{0}

\subsubsection{Preliminaries: the Cauchy problem, and what a difference scheme for advection is}

Problem~5 is not asking whether \(u_t+a u_x=0\) is well-posed. That PDE, on the line, is the model hyperbolic Cauchy problem: the solution is a pure translation, every \(L^p\) norm is conserved, and no boundary condition is needed. The exam asks two discrete questions on top of that object. Part~(i) is an \(\ell_2\) energy identity for the Lax--Wendroff scheme on \(\mathbb{Z}\), which is a stability certificate that does not use Fourier. Part~(ii) changes the domain to a bounded interval, where the Cauchy problem is replaced by an initial-boundary-value problem and Lax--Wendroff, being a three-point stencil, needs a numerical boundary condition at the outflow end.

The official CAM syllabus names this circle as finite-difference methods, stability and convergence, and Lax equivalence. The standard references are LeVeque \cite{leveque07}, Gustafsson--Kreiss--Oliger \cite{gko95}, and Strikwerda \cite{strikwerda04}. The 2026 Spring notes compute the same objects by von Neumann; this paper asks for the energy identity instead. The definitions below are the ones that make the two parts well-posed.

\paragraph{The continuous problem: Cauchy versus initial-boundary-value.}

\begin{definition}[Cauchy problem for constant-coefficient advection]
\label{def:cauchy-adv}
Let \(a\in\mathbb{R}\) be given. The \emph{Cauchy problem} for the advection equation on the real line is: given \(u_0:\mathbb{R}\to\mathbb{R}\), find \(u=u(x,t)\) such that
\begin{equation}
\label{eq:cauchy-adv}
\begin{cases}
u_t + a u_x = 0, & x\in\mathbb{R},\ t>0, \\
u(x,0) = u_0(x), & x\in\mathbb{R}.
\end{cases}
\end{equation}
No boundary condition is imposed: the spatial domain has no boundary. If \(u_0\in L^2(\mathbb{R})\), one seeks a solution in \(C\bigl([0,\infty); L^2(\mathbb{R})\bigr)\). The exam's phrase ``the Cauchy problem imposed on the real line'' is this IVP, discretized on the grid \(\mathbb{Z}\) below.
\end{definition}

The PDE \(u_t+a u_x=0\) is the directional derivative of \(u\) along the vector \((a,1)\) in the \((x,t)\)-plane. Its characteristics are the straight lines \(x-at=\mathrm{const}\).

\begin{proposition}[Solution by characteristics; \(L^2\) conservation]
\label{prop:char}
Assume \(u_0\) is \(C^1\) (or, after density, merely \(L^2\)). The unique solution of \eqref{eq:cauchy-adv} is the travelling wave
\begin{equation}
\label{eq:travelling}
u(x,t)
=
u_0(x-at).
\end{equation}
In particular \(\|u(\,\cdot\,,t)\|_{L^2(\mathbb{R})}=\|u_0\|_{L^2(\mathbb{R})}\) for every \(t\ge 0\), and likewise in every \(L^p\). The problem is well-posed in \(L^2\): the solution operator is an isometry, of norm independent of \(t\).
\end{proposition}

\begin{proof}
Along a characteristic \(x(t)=x_0+at\) one has \(\frac{\mathrm{d}}{\mathrm{d}t}u(x(t),t)=u_t+a u_x=0\), so \(u\) is constant on each such line. The line through \((x,t)\) hits \(t=0\) at \(x-at\), which is \eqref{eq:travelling}. The \(L^2\) identity is translation invariance of Lebesgue measure. Uniqueness for \(C^1\) data is the same characteristic computation run backwards; the \(L^2\) theory is the continuous extension of that isometry.
\end{proof}

Information travels at speed \(a\), and only in one direction. That is the distinction that appears the moment the spatial domain is cut down to an interval.

\begin{definition}[Initial-boundary-value problem; inflow versus outflow]
\label{def:ibvp-adv}
On a bounded interval \((0,1)\), the same PDE must be supplemented by a boundary condition at the \emph{inflow} end, and must not be supplemented by a boundary condition at the \emph{outflow} end. For \(a>0\) the characteristics enter at \(x=0\) and leave at \(x=1\), so the IBVP is
\begin{equation}
\label{eq:ibvp-adv}
\begin{cases}
u_t + a u_x = 0, & 0<x<1,\ t>0, \\
u(x,0) = u_0(x), & 0<x<1, \\
u(0,t) = g(t), & t>0.
\end{cases}
\end{equation}
The exam takes the homogeneous inflow datum \(g\equiv 0\). At \(x=1\) the trace \(u(1,t)\) is determined by the PDE (it equals \(u_0(1-at)\) for \(t<1/a\), and equals \(g(t-1/a)\) thereafter). Imposing an additional boundary condition at \(x=1\) would overdetermine the PDE.
\end{definition}

\begin{center}
\begin{tikzpicture}[>=Latex, scale=1.15]
  \draw[->] (-0.35,0) -- (4.6,0) node[right] {$x$};
  \draw[->] (0,-0.25) -- (0,3.15) node[above] {$t$};
  \draw[thick] (0,0) -- (0,2.85);
  \draw[thick] (3.6,0) -- (3.6,2.85);
  \node[below] at (0,0) {$0$};
  \node[below] at (3.6,0) {$1$};
  \foreach \b in {-0.9,-0.3,0.3,0.9,1.5,2.1,2.7}
    \draw[blue!65!black] ({max(0,\b)}, {max(0,-0.55*\b)}) -- ({min(3.6,\b+3.6)}, {0.55*(min(3.6,\b+3.6)-\b)});
  \node[font=\small, align=center, blue!65!black] at (1.85,2.55) {characteristics \(x-at=\mathrm{const}\)};
  \node[font=\small, rotate=90, anchor=south] at (-0.22,1.35) {inflow};
  \node[font=\small, rotate=90, anchor=south] at (3.82,1.35) {outflow};
  \node[font=\small] at (1.8,-0.42) {Cauchy: the same lines on \(\mathbb{R}\), no walls};
\end{tikzpicture}
\end{center}

A difference scheme does not automatically inherit this distinction. Lax--Wendroff at an interior node reads both neighbours. On \(\mathbb{Z}\) that is harmless. On \(\{0,1,\dots,J\}\) the node adjacent to \(x=1\) asks for a value to the right of the interval, which the PDE does not provide. That missing value is a \emph{numerical} boundary condition (\cref{def:nbc}), and it is the object of part~(ii).

\paragraph{The discrete objects: grid, differences, \(\ell_2\).}

\begin{notation}[Grid and Courant number]
\label{not:grid-p5}
Fix a mesh size \(h>0\) and a time step \(\tau>0\). Write \(x_j=jh\) and \(t_n=n\tau\). A \emph{grid function} \(v=(v_j)\) assigns a value to each node; a time-dependent grid function is written \(u_j^n\), intended to approximate \(u(x_j,t_n)\). For the Cauchy problem the spatial index runs over \(j\in\mathbb{Z}\). For the IBVP on \((0,1)\) one takes \(h=1/J\) and \(j=0,1,\dots,J\), with \(x_0=0\) and \(x_J=1\).

The mesh Courant number (CFL number) is
\[
\nu
\coloneqq
\frac{a\tau}{h}.
\]
It is the number of mesh cells a characteristic travels in one time step. The scheme in the exam is written entirely in terms of \(\nu\).
\end{notation}

\begin{definition}[Forward and backward differences]
\label{def:delta}
For a grid function \(v=(v_j)\), the exam's undivided differences are
\[
\delta_x^+ v_j
\coloneqq
v_{j+1}-v_j,
\qquad
\delta_x^- v_j
\coloneqq
v_j-v_{j-1}.
\]
The corresponding divided (first-order) quotients, which approximate \(\partial_x\), are

\begin{align*}
D^+ v_j
&\coloneqq
\frac{\delta_x^+ v_j}{h}
=
\frac{v_{j+1}-v_j}{h},
\\
D^- v_j
&\coloneqq
\frac{\delta_x^- v_j}{h}
=
\frac{v_j-v_{j-1}}{h},
\\
D^0 v_j
&\coloneqq
\frac{v_{j+1}-v_{j-1}}{2h}
=
\frac{(\delta_x^++\delta_x^-)v_j}{2h}.
\end{align*}

The undivided second difference is
\[
\delta_x^+\delta_x^- v_j
=
\delta_x^-\delta_x^+ v_j
=
v_{j+1}-2v_j+v_{j-1}
=
h^2 D^+D^- v_j,
\]
and \(D^+D^-v_j=(v_{j+1}-2v_j+v_{j-1})/h^2\) approximates \(u_{xx}\) to order \(h^2\). On \(\mathbb{Z}\) one has the shift identity \(\delta_x^- v_j=\delta_x^+ v_{j-1}\).
\end{definition}

Taylor expansion at \(x_j\) gives \(D^\pm v=v_x\pm\tfrac12 h v_{xx}+O(h^2)\) and \(D^0 v=v_x+O(h^2)\), as in the 2026 notes. The energy identity of part~(i) is written in the undivided operators \(\delta_x^\pm\), which differ from \(D^\pm\) only by the factor \(h\); the two languages are interchangeable.

\begin{definition}[Discrete \(\ell_2\) on \(\mathbb{Z}\)]
\label{def:discrete-l2}
For grid functions \(v,w\) on \(\mathbb{Z}\) (real-valued, in the exam),
\[
\|v\|_2^2
\coloneqq
\sum_{j\in\mathbb{Z}}\lvert v_j\rvert^2,
\qquad
\langle v,w\rangle
\coloneqq
\sum_{j\in\mathbb{Z}} v_j w_j,
\]
whenever the sums converge. Equivalently, \(v\in\ell_2(\mathbb{Z})\). This is the discrete analogue of \(L^2(\mathbb{R})\), and it is the space in which the Cauchy problem \eqref{eq:cauchy-adv} is discretized. There is no factor of \(h\) in the exam's definition: \(\|\cdot\|_2\) is the \(\ell_2\) sequence norm, not the trapezoid approximation to \(\|u\|_{L^2}\). For stability the missing \(h\) is a global constant and does not matter; for comparing to \(L^2(\mathbb{R})\) one would use \(h^{1/2}\|v\|_2\).
\end{definition}

On a bounded grid \(\{0,\dots,J\}\) the same formulae are used with the sum restricted to the interior nodes, and boundary terms appear after summation by parts. Part~(i) has no boundary: the line is \(\mathbb{Z}\).

\begin{lemma}[Shift isometries and summation by parts on \(\mathbb{Z}\)]
\label{lem:sbp}
Let \(v,w\in\ell_2(\mathbb{Z})\). Then \(\|\delta_x^+ v\|_2=\|\delta_x^- v\|_2\), and
\begin{equation}
\label{eq:sbp}
\langle \delta_x^+ v,\, w\rangle
=
-
\langle v,\, \delta_x^- w\rangle.
\end{equation}
In particular \(\langle \delta_x^+ v,\, \delta_x^- v\rangle\le\|\delta_x^+ v\|_2^2\), hence
\begin{equation}
\label{eq:dissip-sign}
\|\delta_x^+ v\|_2^2
-
\langle \delta_x^+ v,\, \delta_x^- v\rangle
\ge
0.
\end{equation}
\end{lemma}

\begin{proof}
The map \(v_j\mapsto v_{j-1}\) is an \(\ell_2\) isometry of \(\mathbb{Z}\), and \(\delta_x^- v_j=\delta_x^+ v_{j-1}\), so the two undivided gradients have equal norms. For \eqref{eq:sbp}, expand \(\sum_j(v_{j+1}-v_j)w_j=\sum_j v_{j+1}w_j-\sum_j v_j w_j\) and reindex the first sum (\(j\mapsto j-1\)); both rearrangements are justified by absolute convergence of the Cauchy--Schwarz estimates \(\sum\lvert v_{j+1}w_j\rvert\le\|v\|_2\|w\|_2\). The inequality \eqref{eq:dissip-sign} is Cauchy--Schwarz:
\[
\langle \delta_x^+ v,\, \delta_x^- v\rangle
\le
\|\delta_x^+ v\|_2\,\|\delta_x^- v\|_2
=
\|\delta_x^+ v\|_2^2.
\]
\end{proof}

The remainder in the exam identity is a positive multiple of \eqref{eq:dissip-sign} whenever \(\lvert\nu\rvert\le 1\). That is why an algebraic \(\ell_2\) budget, with no Fourier transform, can prove stability of Lax--Wendroff on the line.

\paragraph{Difference schemes: consistency, stability, convergence.}

\begin{definition}[Linear one-step scheme]
\label{def:scheme}
A linear one-step finite-difference scheme on \(\mathbb{Z}\) is a marching rule
\[
U^{n+1}
=
S_{h,\tau} U^n,
\]
where \(S_{h,\tau}\) is a bounded linear operator on \(\ell_2(\mathbb{Z})\) (typically a finite stencil, i.e.\ a Laurent polynomial in the shift). The scheme is \emph{explicit} if \(U^{n+1}\) is a finite linear combination of already-known values \(U^n_{\,j+k}\); it is \emph{implicit} if computing \(U^{n+1}\) requires solving a linear system that couples several nodes at time \(t_{n+1}\). Lax--Wendroff is explicit and three-point.
\end{definition}

\begin{definition}[Local truncation error / consistency]
\label{def:truncation-p5}
Write the scheme in residual form \(\mathcal{L}_{h,\tau}U=0\). Let \(u\) be a smooth solution of \(u_t+a u_x=0\), and write \(u_j^n\coloneqq u(x_j,t_n)\) for its samples. The \emph{local truncation error} is the residual of those samples:
\[
\mathcal{T}_j^{n+1}
\coloneqq
\bigl(\mathcal{L}_{h,\tau} u\bigr)_j^{n+1}.
\]
The scheme is \emph{consistent} if \(\mathcal{T}_j^n\to 0\) as \(h,\tau\to 0\), at every fixed \((x,t)\), for every smooth solution \(u\). The order is the largest pair \((p,q)\) with \(\mathcal{T}=O(\tau^p+h^q)\). Consistency is a statement about \(u\), not about the computed grid function \(U\).
\end{definition}

\begin{definition}[Lax--Richtmyer stability]
\label{def:stability-p5}
The scheme \(U^{n+1}=S_{h,\tau}U^n\) is \emph{stable} on a time interval \([0,T]\), in the discrete norm \(\|\cdot\|_h\), if there exist \(C,\alpha\), independent of \(h\) and \(\tau\) (for all sufficiently small steps, possibly subject to a restriction such as \(\lvert\nu\rvert\le 1\)), such that
\[
\|U^n\|_h
\le
C\,e^{\alpha t_n}\,\|U^0\|_h,
\qquad
0\le t_n\le T.
\]
For advection, whose \(L^2\) solution operator is an isometry (\cref{prop:char}), the practical target is \(\|U^n\|_2\le\|U^0\|_2\), i.e.\ \(C=1\) and \(\alpha=0\). An identity of the form \(\|U^{n+1}\|_2^2=\|U^n\|_2^2-\text{(nonnegative)}\) is a stability proof in this sense.
\end{definition}

\begin{theorem}[Lax equivalence]
\label{thm:lax-eq-p5}
For a linear well-posed evolutionary IVP and a consistent finite-difference scheme, convergence of the grid function to the PDE solution (in the discrete norm, as \(h,\tau\to 0\) with \(t_n\le T\)) holds if and only if the scheme is stable \cite{laxrichtmyer56,leveque07,gko95}.
\end{theorem}

The exam does not ask one to quote \cref{thm:lax-eq-p5}. It asks one to prove the \(\ell_2\) bound that makes the ``if'' direction applicable. Von Neumann's criterion (Fourier substitution \(u_j^n=g(\xi)^n e^{\mathrm{i} j\xi}\), requiring \(\lvert g(\xi)\rvert\le 1\) for advection) is equivalent to that bound for constant-coefficient schemes on \(\mathbb{Z}\), and is the machine used on the 2026 paper. Part~(i) is the same bound, obtained by inner products.

\paragraph{Lax--Wendroff.}

\begin{definition}[Lax--Wendroff scheme]
\label{def:lw}
The Lax--Wendroff scheme for \(u_t+a u_x=0\) is the explicit three-point update
\begin{equation}
\label{eq:lw}
u_j^{n+1}
=
-\tfrac12\nu(1-\nu)\,u_{j+1}^n
+
(1-\nu^2)\,u_j^n
+
\tfrac12\nu(1+\nu)\,u_{j-1}^n,
\end{equation}
with \(\nu=a\tau/h\). In undivided differences (\cref{def:delta}),
\begin{equation}
\label{eq:lw-delta}
u^{n+1}
=
u^n
-
\frac{\nu}{2}(\delta_x^++\delta_x^-)u^n
+
\frac{\nu^2}{2}\,\delta_x^+\delta_x^- u^n.
\end{equation}
The two forms are identical: \(\delta_x^++\delta_x^-\) is the centered undivided first difference, and \(\delta_x^+\delta_x^-\) is the undivided second difference.
\end{definition}

The scheme is the unique three-point explicit method that matches the Taylor expansion of the travelling wave through second order in \(\tau\). That is its definition, not an independent accuracy check.

\begin{proposition}[Derivation of Lax--Wendroff]
\label{prop:lw-derive}
Let \(u\) be a smooth solution of \(u_t+a u_x=0\). Then \(u_{tt}=a^2 u_{xx}\), and
\[
u(x,t+\tau)
=
u
-
a\tau\, u_x
+
\frac{a^2\tau^2}{2}\,u_{xx}
+
O(\tau^3),
\]
the expansions being at \((x,t)\). Replacing \(u_x\) by the centered quotient \(D^0\) and \(u_{xx}\) by \(D^+D^-\), both at time \(t_n\), and dropping \(O(\tau^3+h^2\tau+\tau h^2)\) produces \eqref{eq:lw}.
\end{proposition}

\begin{proof}
Taylor in time:
\[
u(x,t+\tau)
=
u+\tau u_t+\tfrac12\tau^2 u_{tt}+O(\tau^3).
\]
The PDE gives \(u_t=-a u_x\). Differentiating it in \(t\) and in \(x\) and comparing mixed derivatives gives \(u_{tt}=-a u_{xt}=-a\partial_x(u_t)=a^2 u_{xx}\). Substitute. On the grid, \(a\tau=\nu h\), and the standard centered replacements
\[
u_x
=
D^0 u+O(h^2),
\qquad
u_{xx}
=
D^+D^- u+O(h^2)
\]
turn the right-hand side into
\[
u_j^n
-
\frac{\nu}{2}\bigl(u_{j+1}^n-u_{j-1}^n\bigr)
+
\frac{\nu^2}{2}\bigl(u_{j+1}^n-2u_j^n+u_{j-1}^n\bigr)
+
O(\tau^3+\tau h^2).
\]
Collecting coefficients of \(u_{j+1}^n\), \(u_j^n\), and \(u_{j-1}^n\) is \eqref{eq:lw}. The operator form \eqref{eq:lw-delta} is the same collection written with \(\delta_x^\pm\).
\end{proof}

\begin{proposition}[Truncation of Lax--Wendroff]
\label{prop:lw-trunc}
On smooth solutions of \(u_t+a u_x=0\), the local truncation error of \eqref{eq:lw} is \(O(\tau^2+h^2)\). The scheme is consistent of order \(2\) in time and space. The leading terms may be written
\[
\mathcal{T}
=
-\frac{\tau^2}{6}u_{ttt}
-
\frac{a h^2}{6}u_{xxx}
+
\frac{a^2\tau h^2}{24}(\cdots)
+
O(\tau^3+h^4),
\]
after using \(u_{tt}=a^2 u_{xx}\) to replace mixed derivatives; the precise cubic polynomial is not needed on this paper.
\end{proposition}

The derivation already shows that every term through \(O(\tau^2)\) in time and \(O(h^2)\) in space has been matched, so the residual starts at the next order. Combined with stability under \(\lvert\nu\rvert\le 1\) (\cref{prop:lw-vn,prop:lw-energy-sign}), Lax equivalence yields second-order convergence on the Cauchy problem.

\begin{definition}[CFL]
\label{def:cfl-p5}
An explicit scheme's \emph{numerical domain of dependence} at \((x_j,t_{n+1})\) is the set of nodes at time \(t_n\) that enter the update of \(u_j^{n+1}\). Courant--Friedrichs--Lewy: a necessary condition for convergence is that this set contain the characteristic foot \(x_j-a\tau\). For any explicit three-point stencil the numerical domain is \(\{x_{j-1},x_j,x_{j+1}\}\), hence the necessary restriction is
\[
\lvert a\tau\rvert
\le
h
\qquad\text{i.e.}\qquad
\lvert\nu\rvert\le 1.
\]
\end{definition}

CFL is not a stability theorem. Explicit centered advection (\(u_j^{n+1}=u_j^n-\frac{\nu}{2}(u_{j+1}^n-u_{j-1}^n)\)) satisfies \(\lvert\nu\rvert\le 1\) as a domain-of-dependence restriction and is nevertheless unconditionally unstable by von Neumann. For Lax--Wendroff the CFL restriction and the \(\ell_2\) restriction happen to coincide.

\begin{proposition}[von Neumann symbol of Lax--Wendroff]
\label{prop:lw-vn}
Substituting \(u_j^n=g^n e^{\mathrm{i} j\xi}\) into \eqref{eq:lw} produces
\[
g(\xi)
=
1-\nu^2\bigl(1-\cos\xi\bigr)
-
\mathrm{i}\,\nu\sin\xi,
\]
and
\begin{equation}
\label{eq:lw-gmod}
\lvert g(\xi)\rvert^2
=
1-4\nu^2(1-\nu^2)\sin^4(\xi/2).
\end{equation}
Thus \(\lvert g(\xi)\rvert\le 1\) for all \(\xi\) if and only if \(\lvert\nu\rvert\le 1\). Under that restriction the scheme is Lax--Richtmyer stable in \(\ell_2(\mathbb{Z})\). For \(\lvert\nu\rvert<1\) one has \(\lvert g\rvert<1\) except at \(\xi\in 2\pi\mathbb{Z}\): the scheme is dissipative on high modes.
\end{proposition}

\begin{proof}
The symbol of \(\delta_x^+\) is \(e^{\mathrm{i}\xi}-1\), of \(\delta_x^-\) is \(1-e^{-\mathrm{i}\xi}\), and of \(\delta_x^++\delta_x^-\) is \(2\mathrm{i}\sin\xi\). Substituting \eqref{eq:lw-delta} gives
\[
g
=
1
-
\frac{\nu}{2}\cdot 2\mathrm{i}\sin\xi
+
\frac{\nu^2}{2}\cdot\bigl(e^{\mathrm{i}\xi}-2+e^{-\mathrm{i}\xi}\bigr)
=
1-\mathrm{i}\nu\sin\xi-\nu^2(1-\cos\xi).
\]
Then \(\lvert g\rvert^2=(1-\nu^2(1-\cos\xi))^2+\nu^2\sin^2\xi\). Expanding and using \(1-\cos\xi=2\sin^2(\xi/2)\), \(\sin\xi=2\sin(\xi/2)\cos(\xi/2)\) produces \eqref{eq:lw-gmod}. The factor \(4\nu^2(1-\nu^2)\sin^4(\xi/2)\) is nonnegative for all \(\xi\) precisely when \(\nu^2\in[0,1]\).
\end{proof}

The same factor \(\nu^2(1-\nu^2)\) sits in front of the discrete gradient remainder in part~(i). Fourier and the energy identity are two writings of one dissipation.

\begin{proposition}[Sign of the exam remainder]
\label{prop:lw-energy-sign}
If \(\lvert\nu\rvert\le 1\) and \(u^n\in\ell_2(\mathbb{Z})\), then
\[
\frac12\nu^2(1-\nu^2)
\bigl(
\|\delta_x^+ u^n\|_2^2
-
\langle \delta_x^+ u^n,\,\delta_x^- u^n\rangle
\bigr)
\ge
0.
\]
Consequently any identity of the form
\[
\|u^{n+1}\|_2^2
=
\|u^n\|_2^2
-
\frac12\nu^2(1-\nu^2)
\bigl(
\|\delta_x^+ u^n\|_2^2
-
\langle \delta_x^+ u^n,\,\delta_x^- u^n\rangle
\bigr)
\]
implies \(\|u^{n+1}\|_2\le\|u^n\|_2\), hence Lax--Richtmyer stability on the Cauchy grid. The identity itself is the algebraic content of part~(i): expand \(\|u^{n+1}\|_2^2\) from \eqref{eq:lw-delta}, and reduce the cross terms by \cref{lem:sbp}.
\end{proposition}

\paragraph{Numerical boundary conditions.}

\begin{definition}[Physical vs.\ numerical boundary conditions]
\label{def:nbc}
On a bounded grid \(j=0,1,\dots,J\),
\begin{enumerate}[label=\textup{(\roman*)}]
    \item a \emph{physical} (or PDE) boundary condition is a prescription of boundary data that the continuous IBVP \eqref{eq:ibvp-adv} actually requires;
    \item a \emph{numerical} boundary condition is an additional closure, needed by the stencil but not by the PDE, that expresses an unknown near-boundary value in terms of interior (and physical-boundary) values.
\end{enumerate}
For \(a>0\), the physical condition is at inflow: \(u_0^n=g(t_n)\), and the exam takes \(g=0\). Lax--Wendroff at \(j=J-1\) (or at \(j=J\), depending on whether one treats \(x_J=1\) as an unknown) reads \(u_J^n\) and \(u_{J+1}^n\). The value \(u_{J+1}^n\) lies outside \((0,1)\) and is not supplied by \eqref{eq:ibvp-adv}. Closing it is a numerical boundary condition at outflow.
\end{definition}

A scheme that is stable on \(\mathbb{Z}\) can be destabilized by a bad numerical boundary condition. The precise theory is Gustafsson--Kreiss--Sundstr\"om (GKS): one freezes coefficients at the boundary, substitutes modes \(z^n\kappa^j\), and forbids incoming \(z\)-roots with \(\lvert z\rvert>1\) that the interior scheme does not already control \cite{gko95}. The exam asks for a \emph{simple} condition that is known to pass, not for the GKS test.

\begin{construction}[A simple stable outflow condition, \(a>0\)]
\label{con:nbc-simple}
Keep the homogeneous physical condition \(u_0^n=0\). At the right end, do not apply Lax--Wendroff. Close the last unknown by the first-order upwind (backward Euler in space, forward Euler in time) scheme
\begin{equation}
\label{eq:nbc-upwind}
\frac{u_J^{n+1}-u_J^n}{\tau}
+
a\,\frac{u_J^n-u_{J-1}^n}{h}
=
0,
\qquad\text{i.e.}\qquad
u_J^{n+1}
=
(1-\nu)\,u_J^n
+
\nu\,u_{J-1}^n.
\end{equation}
This uses only values inside \([0,1]\), is consistent of order \(O(\tau+h)\) with the PDE, and is stable for \(0\le\nu\le 1\) (it is the standard explicit upwind scheme, whose amplification factor is \(g=(1-\nu)+\nu e^{-\mathrm{i}\xi}\) with \(\lvert g\rvert\le 1\) iff \(0\le\nu\le 1\)). The interior Lax--Wendroff updates occupy \(j=1,\dots,J-1\).

An equally standard alternative is first-order spatial extrapolation of a ghost node,
\[
u_{J+1}^n
=
u_J^n
\qquad\text{or}\qquad
u_{J+1}^n
=
2u_J^n-u_{J-1}^n,
\]
followed by Lax--Wendroff at \(j=J\). The constant extrapolation \(u_{J+1}^n=u_J^n\) is the simpler of the two; under \(0\le\nu\le 1\) it is GKS-stable for this outflow problem \cite{gko95}. Either closure is an acceptable answer to part~(ii). What is \emph{not} acceptable is to impose a physical Dirichlet condition at \(x=1\) (the PDE does not allow it), or to use a downwind stencil that looks to the right of \(x=1\) without defining the ghost.
\end{construction}

The order drop at a single boundary node does not destroy the global \(\ell_2\) rate on a fixed interval: the boundary is codimension one. For the exam, stability is the requirement, not preservation of second order.

\begin{examnote}[What the two parts are]
\label{examnote:p5}
\begin{enumerate}[label=\textup{(\roman*)}]
    \item Cauchy problem: \cref{def:cauchy-adv}, discrete space \cref{def:discrete-l2}. The scheme is \cref{def:lw} in the form \eqref{eq:lw-delta}. Expand \(\|u^{n+1}\|_2^2\), reduce shifts by \cref{lem:sbp}, and identify the remainder as the displayed combination of \(\|\delta_x^+\|_2^2\) and \(\langle\delta_x^+,\delta_x^-\rangle\). The sign \eqref{eq:dissip-sign} plus \(\lvert\nu\rvert\le 1\) is \(\ell_2\) stability, which is the reason the identity is asked.
    \item IBVP: \cref{def:ibvp-adv}. Physical BC at \(x=0\); numerical BC at \(x=1\), \cref{def:nbc,con:nbc-simple}. Quote upwind \eqref{eq:nbc-upwind} or constant extrapolation, and the restriction \(0\le\nu\le 1\).
\end{enumerate}
\end{examnote}

\begin{remark}[Relation to 2026 Spring Problems~4--5]
\label{rem:p5-vs-2026}
The 2026 P4 notes develop von Neumann, CFL, truncation, and Lax equivalence for three \emph{implicit} advection schemes. The objects in \cref{def:truncation-p5,def:stability-p5,thm:lax-eq-p5,def:cfl-p5} are the same objects, specialized to explicit Lax--Wendroff. The 2026 P5 notes develop a discrete maximum principle (\(\ell_\infty\), Dirichlet, variable coefficients). That machine does not apply to this paper: Lax--Wendroff is not monotone for all \(\nu\in(0,1)\), and part~(i) is an \(\ell_2\) identity on the line.
\end{remark}


\noindent\textbf{Remark: You can choose either 6 or 7. The points will be decided as \(\max(6,7)\).}

\subsection{Problem 6. [20 points]}
Consider a harmonic oscillator with a cubic damping term

\[
y'' + y + \epsilon (y')^{3} = 0, \qquad y(0) = 1,\quad y'(0) = 0,
\]

where \(y = y(t)\), \(t \ge 0\), \(\epsilon > 0\).
\begin{enumerate}[label=(\roman*)]
    \item For small \(\epsilon\), use the multiple-scale method to study the behavior of \(y(t)\) for large \(t\), i.e., construct a proper asymptotic solution.
    \item Make a conclusion on the validity of your asymptotic solution. Briefly justify your conclusion.
\end{enumerate}

\subsection{Problem 7. [20 points]}
Let \(K\) be a nonempty closed convex set in \(\mathbb{R}^{n}\). For any \(z \in \mathbb{R}^{n}\), define \(\Pi_{K}(z)\) as the optimal solution of the problem \(\min \tfrac12 \|y - z\|\), s.t.\ \(y \in K\). Prove the following statements.
\begin{enumerate}[label=(\roman*)]
    \item Show that \(\Pi_{K}(z)\) satisfies
    \[
    \langle z - \Pi_{K}(z), d - \Pi_{K}(z) \rangle \le 0, \qquad \forall d \in K.
    \]
    \item Define
    \[
    \Theta(z) := \tfrac12 \|z - \Pi_{K}(z)\|^{2}.
    \]
    Prove that \(\Theta(\cdot)\) is continuous differentiable convex function and
    \[
    \nabla \Theta(z) = z - \Pi_{K}(z).
    \]
\end{enumerate}

\begin{thebibliography}{99}
\bibitem[Qingshu(11th)]{qingshu11}
清疏.
\newblock \emph{大学生数学竞赛班第十一届讲义}.
\newblock Euler--Maclaurin: \S8.5, Theorem~8.8 / (8.42), Corollary~8.3 / (8.44). Periodic Bernoulli functions \(b_{n}\): Definition~8.1, Proposition~8.4, Corollary~8.2. Local slice: \texttt{qe-review/applied-math/reference/qingshu-euler-maclaurin/}.
\bibitem[ConteDeBoor(2000)]{conte00}
S.~D. Conte and C.~de~Boor.
\newblock \emph{Elementary Numerical Analysis: An Algorithmic Approach}, 3rd~ed.
\newblock McGraw--Hill, 1980; reprint SIAM, 2017.
\newblock Official CAM syllabus. Newton iteration, order of convergence, and multiple roots. Trigonometric interpolation / Fourier: the aliasing argument of Problem~1(ii).
\bibitem[LeVeque(2007)]{leveque07}
R.~J. LeVeque.
\newblock \emph{Finite Difference Methods for Ordinary and Partial Differential Equations}.
\newblock SIAM, 2007.
\newblock Difference quotients and truncation: Ch.~1. Advection, CFL, von Neumann, Lax equivalence: Chs.~8--10. Lax--Wendroff: Ch.~10.
\bibitem[Strikwerda(2004)]{strikwerda04}
J.~C. Strikwerda.
\newblock \emph{Finite Difference Schemes and Partial Differential Equations}.
\newblock SIAM, 2004.
\newblock Hyperbolic Fourier analysis: Chs.~5--6.
\bibitem[GKO(1995)]{gko95}
B.~Gustafsson, H.-O.~Kreiss, and J.~Oliger.
\newblock \emph{Time Dependent Problems and Difference Methods}.
\newblock Wiley, 1995.
\newblock Official CAM syllabus item~5. Fourier multipliers; GKS theory of numerical boundary conditions.
\bibitem[LaxRichtmyer(1956)]{laxrichtmyer56}
P.~D. Lax and R.~D. Richtmyer.
\newblock Survey of the stability of linear finite difference equations.
\newblock \emph{Comm.\ Pure Appl.\ Math.} 9 (1956), 267--293.
\newblock Lax equivalence: a consistent approximation to a linear well-posed IVP converges if and only if it is stable.
\end{thebibliography}

\end{document}
